% Generated mechanically from the frozen moduli-curves.tex target.
% Do not edit this generated file; edit the bound locale source instead.

% OK, start here.
%

\setcounter{chapter}{108}
\chapter{曲线模空间}


\phantomsection
\label{moduli-curves-section-phantom}





\section{引言}
\label{moduli-curves-section-introduction}

\noindent
本章讨论若干熟知的曲线模叠。一个参考文献是
Deligne 与 Mumford 的著名论文，见 \cite{DM}。




\section{约定与语言滥用}
\label{moduli-curves-section-conventions}

\noindent
我们继续采用《叠的性质》中的章节
\ref{stacks-properties-section-conventions} 所引入的约定和语言滥用。
除非另有说明，基概形均为 $\Spec(\mathbf{Z})$。






\section{曲线叠}
\label{moduli-curves-section-stack-curves}

\noindent
本节是《Quot》中的章节 \ref{quot-section-curves} 的续篇。
设 $\Curvesstack$ 为这样一个叠：它在概形 $S$ 上的截面范畴是
$S$ 上曲线族的范畴。需要特别记住，所谓
{\it 曲线族}，是指态射 $f : X \to S$，其中 $X$
是代数空间（！），而 $f$ 平坦、真、有限表示且相对维数
$\leq 1$。我们已经知道 $\Curvesstack$ 是
$\mathbf{Z}$ 上的代数叠，见《Quot》中的定理
\ref{quot-theorem-curves-algebraic}。
如果在叠的定义中不允许代数空间，则该定理并不成立。

\medskip\noindent
基变换通常用下标表示，但对 $\Curvesstack$ 不能采用这种记号，
因为 $\Curvesstack_S$ 已用来表示 $S$ 上的纤维范畴。
因此，我们在《Quot》中的注
\ref{quot-remark-curves-base-change} 里用 $B\text{-}\Curvesstack$
表示到代数空间 $B$ 的基变换
$$
B\text{-}\Curvesstack = \Curvesstack \times B
$$
右侧的积取在终对象上，即取在 $\Spec(\mathbf{Z})$ 上。
左侧对象是这样一个叠：它分类 $B$ 上概形范畴中的曲线族。
特别地，若 $k$ 是域，则
$$
k\text{-}\Curvesstack = \Curvesstack \times \Spec(k)
$$
是分类 $k$ 上概形范畴中的曲线族的模叠。
在继续之前，先作如下相容性检验。

\begin{lemma}
\label{moduli-curves-lemma-extend-curves-to-spaces}
设 $T \to B$ 是代数空间的态射。范畴
$$
\Mor_B(T, B\text{-}\Curvesstack) = \Mor(T, \Curvesstack)
$$
就是 $T$ 上曲线族的范畴。
\end{lemma}

\begin{proof}
$T$ 上的曲线族是代数空间的态射 $f : X \to T$，它平坦、真、
有限表示且相对维数 $\leq 1$（《空间的态射》中的定义
\ref{spaces-morphisms-definition-relative-dimension}）。
这与《Quot》中的情形 \ref{quot-situation-curves} 的定义完全相同，
不同之处仅在于允许基 $T$ 为代数空间。
代数叠和代数空间的默认基范畴是概形范畴，因此本引理并不直接
由定义得出。不过，我们仍建议读者略过此证明。

\medskip\noindent
由上面给出的 $B\text{-}\Curvesstack$ 的积描述，只需证明绝对情形。
取概形 $U$ 以及满的 \'etale 态射 $p : U \to T$。
令 $R = U \times_T U$，其投影为 $s, t : R \to U$。

\medskip\noindent
设 $v : T \to \Curvesstack$ 是态射。则 $v \circ p$
对应于曲线族 $X_U \to U$。典范 $2$-态射
$v \circ p \circ t \to v \circ p \circ s$
是一个同构
$\varphi : X_U \times_{U, s} R \to X_U \times_{U, t} R$。
该同构在 $R \times_{s, t} R$ 上满足余圈条件。
由《自举》中的引理 \ref{bootstrap-lemma-descend-algebraic-space}，
得到代数空间的态射 $X \to T$；它在 $U$ 上的拉回等于 $X_U$，
并与 $\varphi$ 相容。由于 $\{U \to T\}$ 是 \'etale 覆盖，
《空间的下降》中的引理
\ref{spaces-descent-lemma-descending-property-flat}、
\ref{spaces-descent-lemma-descending-property-proper} 和
\ref{spaces-descent-lemma-descending-property-finite-presentation}
表明 $X \to T$ 平坦、真且有限表示。
此外，$X \to T$ 的相对维数 $\leq 1$，因为这是 \'etale 局部性质。
故 $X \to T$ 是 $T$ 上的曲线族。

\medskip\noindent
反之，设 $X \to T$ 是曲线族。则基变换 $X_U$ 确定态射
$w : U \to \Curvesstack$，而典范同构
$X_U \times_{U, s} R \to X_U \times_{U, t} R$
确定满足余圈条件的 $2$-箭头 $w \circ s \to w \circ t$。
% P12-ERR-0229: see ../control/ERRATA_CANDIDATES.jsonl
因此，由商 $[U/R]$ 的泛性质得到态射
$v : T = [U/R] \to \Curvesstack$；见《空间中的群胚》中的引理
\ref{spaces-groupoids-lemma-quotient-stack-2-coequalizer}。
% P12-ERR-0230: see ../control/ERRATA_CANDIDATES.jsonl
（事实上，在此情形下，更容易回到引入语言滥用之前，直接构造函子
$\Sch/T \to \Curvesstack$；该函子“就是”态射
$T \to \Curvesstack$。）

\medskip\noindent
我们略去如下验证：上述构造可扩张到对象之间的态射，且彼此拟逆。
\end{proof}













\section{极化曲线叠}
\label{moduli-curves-section-polarized-curves}

\noindent
本节展开《Quot》中的注
\ref{quot-remark-alternative-approach-curves} 所讨论的部分内容。
考虑 $2$-纤维积
$$
\xymatrix{
\Curvesstack \times_{\Spacesstack'_{fp, flat, proper}}
\Polarizedstack \ar[r] \ar[d] &
\Polarizedstack \ar[d] \\
\Curvesstack \ar[r] &
\Spacesstack'_{fp, flat, proper}
}
$$
我们将此 $2$-纤维积记为
$$
\textit{PolarizedCurves} =
\Curvesstack
\times_{\Spacesstack'_{fp, flat, proper}}
\Polarizedstack
$$
该纤维积参数化极化曲线，即赋有相对丰沛可逆层的曲线族。
更确切地，$\textit{PolarizedCurves}$ 的对象是二元组
$(X \to S, \mathcal{L})$，其中
\begin{enumerate}
\item $X \to S$ 是真、平坦、有限表示且相对维数
$\leq 1$ 的概形态射；以及
\item $\mathcal{L}$ 是在 $X/S$ 上相对丰沛的可逆
$\mathcal{O}_X$-模。
\end{enumerate}
$\textit{PolarizedCurves}$ 的对象之间的态射
$(X' \to S', \mathcal{L}') \to (X \to S, \mathcal{L})$
由三元组 $(f, g, \varphi)$ 给出，其中
$f : X' \to X$ 和 $g : S' \to S$ 是概形态射，并组成交换图
$$
\xymatrix{
X' \ar[d] \ar[r]_f & X \ar[d] \\
S' \ar[r]^g & S
}
$$
该图诱导同构 $X' \to S' \times_S X$；换言之，该图是 Cartesian 的，
且 $\varphi : f^*\mathcal{L} \to \mathcal{L}'$ 是同构。
复合按显然的方式定义。

\begin{lemma}
\label{moduli-curves-lemma-polarized-curves-in-polarized}
态射
$\textit{PolarizedCurves} \to
\Polarizedstack$ 是开闭浸入。
\end{lemma}

\begin{proof}
这是因为 $1$-态射
$\Curvesstack \to \Spacesstack'_{fp, flat, proper}$
可表为开闭浸入；见《Quot》中的引理
\ref{quot-lemma-curves-open-and-closed-in-spaces}。
\end{proof}

\begin{lemma}
\label{moduli-curves-lemma-polarized-curves-over-curves}
态射
$\textit{PolarizedCurves} \to \Curvesstack$
光滑且满。
\end{lemma}

\begin{proof}
满性。给定域 $k$ 以及 $k$ 上维数 $\leq 1$ 的真代数空间
$X$，即 $\Curvesstack$ 在 $k$ 上的对象。
由《域上的空间》中的引理
\ref{spaces-over-fields-lemma-codim-1-point-in-schematic-locus}，
代数空间 $X$ 是概形。因此 $X$ 是 $k$ 上维数 $\leq 1$ 的真概形。
由《簇》中的引理 \ref{varieties-lemma-dim-1-proper-projective}，
% P12-ERR-0231: see ../control/ERRATA_CANDIDATES.jsonl
可见 $X$ 在 $\kappa$ 上是 H-射影的。
特别地，$X$ 上存在丰沛可逆 $\mathcal{O}_X$-模
$\mathcal{L}$。于是 $(X, \mathcal{L})$ 是
$\textit{PolarizedCurves}$ 在 $k$ 上的对象，并映到 $X$。

\medskip\noindent
光滑性。设 $X \to S$ 是 $\Curvesstack$ 的对象，即态射
$S \to \Curvesstack$。显然
$$
\textit{PolarizedCurves}
\times_{\Curvesstack} S
\subset \Picardstack_{X/S}
$$
是由如下对象 $(T/S, \mathcal{L}/X_T)$ 构成的子叠：
$\mathcal{L}$ 在 $X_T/T$ 上丰沛。由《空间的下降》中的引理
\ref{spaces-descent-lemma-ample-in-neighbourhood}，这是开子叠。
又由《模叠》中的引理 \ref{moduli-lemma-pic-curves-smooth}，
$\Picardstack_{X/S} \to S$ 光滑，结论遂得。
\end{proof}

\begin{lemma}
\label{moduli-curves-lemma-etale-locally-scheme}
设 $X \to S$ 是曲线族。则存在 \'etale 覆盖
$\{S_i \to S\}$，使得 $X_i = X \times_S S_i$ 是概形。
甚至可以假设 $X_i$ 在 $S_i$ 上是 H-射影的。
\end{lemma}

\begin{proof}
这是引理 \ref{moduli-curves-lemma-polarized-curves-over-curves} 的直接推论。
% P12-ERR-0232: see ../control/ERRATA_CANDIDATES.jsonl
具体而言，展开定义，该引理给出一个满光滑态射 $S' \to S$，
使得 $X' = X \times_S S'$ 赋有在 $X'/S'$ 上丰沛的可逆
$\mathcal{O}_{X'}$-模 $\mathcal{L}'$。
于是可用一个 \'etale 覆盖 $\{S_i \to S\}$ 加细光滑覆盖
$\{S' \to S\}$；见《态射进阶》中的引理
\ref{more-morphisms-lemma-etale-dominates-smooth}。
用适当开覆盖替换 $S_i$ 后，可以假设 $X_i \to S_i$ 是
H-射影的；见《态射》中的引理
\ref{morphisms-lemma-proper-ample-locally-projective} 和
\ref{morphisms-lemma-characterize-locally-projective}
（《态射进阶》中的章节
\ref{more-morphisms-section-projective} 也对此作了详细讨论）。
\end{proof}




\section{曲线叠的性质}
\label{moduli-curves-section-properties}

\noindent
下面的引理对曲面的模空间并不成立；见注
\ref{moduli-curves-remark-boundedness-aut-does-not-work-surfaces}。

\begin{lemma}
\label{moduli-curves-lemma-curves-diagonal-separated-fp}
$\Curvesstack$ 的对角态射分离且有限表示。
\end{lemma}

\begin{proof}
回忆 $\Curvesstack$ 是保持极限的代数叠；见《Quot》中的引理
\ref{quot-lemma-curves-limits}。
由《叠的极限》中的引理
\ref{stacks-limits-lemma-limit-preserving-diagonal}，这蕴含
% P12-ERR-0233: see ../control/ERRATA_CANDIDATES.jsonl
$\Delta : \Polarizedstack \to \Polarizedstack \times \Polarizedstack$
保持极限。因此，由《叠的极限》中的命题
\ref{stacks-limits-proposition-characterize-locally-finite-presentation}，
$\Delta$ 局部有限表示。

\medskip\noindent
下面证明 $\Delta$ 分离。为此，只需说明：给定概形 $U$ 以及
$\Curvesstack$ 在 $U$ 上的两个对象 $Y \to U$ 和 $X \to U$，
代数空间
$$
\mathit{Isom}_U(Y, X)
$$
分离。
% P12-ERR-0234: see ../control/ERRATA_CANDIDATES.jsonl
这已经由《模叠》中的引理 \ref{moduli-lemma-Mor-s-lfp} 和
\ref{moduli-lemma-Isom-in-Mor} 得出：目标是分离代数空间。

\medskip\noindent
为完成证明，我们说明 $\Delta$ 拟紧。由于 $\Delta$
可由代数空间表示，只需检验 $\Delta$ 沿满光滑态射
$U \to \Curvesstack \times \Curvesstack$ 的基变换拟紧
（例如见《叠的性质》中的引理
\ref{stacks-properties-lemma-check-property-covering}）。
取 $U = \coprod U_i$ 为仿射开集的不交并，并取满光滑态射
$$
U \longrightarrow
\textit{PolarizedCurves} \times \textit{PolarizedCurves}
$$
由于引理 \ref{moduli-curves-lemma-polarized-curves-over-curves} 表明
$\textit{PolarizedCurves} \to \Curvesstack$ 满且光滑，
$U \to \Curvesstack \times \Curvesstack$ 也满且光滑。
又因为 $\textit{PolarizedCurves}$ 保持极限
（由《Artin 公理》中的引理
\ref{artin-lemma-fibre-product-limit-preserving} 以及《Quot》中的引理
\ref{quot-lemma-curves-limits}、
\ref{quot-lemma-polarized-limits} 和
\ref{quot-lemma-spaces-limits}），可见
$\textit{PolarizedCurves} \to \Spec(\mathbf{Z})$ 局部有限表示，
因而 $U_i \to \Spec(\mathbf{Z})$ 局部有限表示
（《叠的极限》中的命题
\ref{stacks-limits-proposition-characterize-locally-finite-presentation}
以及《叠的态射》中的引理
\ref{stacks-morphisms-lemma-composition-finite-presentation} 和
\ref{stacks-morphisms-lemma-smooth-locally-finite-presentation}）。
特别地，$U_i$ 是 Noether 仿射概形。这就化归到下一段所讨论的情形。

\medskip\noindent
在本段中，给定 Noether 仿射概形 $U$ 以及
$\textit{PolarizedCurves}$ 在 $U$ 上的两个对象
$(Y, \mathcal{N})$ 和 $(X, \mathcal{L})$，我们证明代数空间
$$
\mathit{Isom}_U(Y, X)
$$
拟紧。由于 $U$ 的连通分支既开又闭，可以用这些连通分支替换
$U$。因此可以并确实假设 $U$ 连通。取一点 $u \in U$。
令 $Q$、$P$ 为这些族的 Hilbert 多项式，即
$$
Q(n) = \chi(Y_u, \mathcal{N}_u^{\otimes n})
\quad\text{且}\quad
P(n) = \chi(X_u, \mathcal{L}_u^{\otimes n})
$$
见《簇》中的引理
\ref{varieties-lemma-numerical-polynomial-from-euler}。
由于 $U$ 连通，且函数
$u \mapsto \chi(Y_u, \mathcal{N}_u^{\otimes n})$ 和
$u \mapsto \chi(X_u, \mathcal{L}_u^{\otimes n})$
局部常值（见《概形的导出范畴》中的引理
\ref{perfect-lemma-chi-locally-constant-geometric}），
可见在 $U$ 的每一点都得到同一 Hilbert 多项式。令
$$
\mathcal{M} = \text{pr}_1^*\mathcal{N}
\otimes_{\mathcal{O}_{Y \times_U X}} \text{pr}_2^*\mathcal{L}
$$
为 $Y \times_U X$ 上的模。给定 $U$ 上某概形 $T$ 以及
% P12-ERR-0235: see ../control/ERRATA_CANDIDATES.jsonl
$(f, \varphi) \in \mathit{Isom}_U(Y, X)(T)$，则对每个 $t \in T$，
有
\begin{align*}
\chi(Y_t, (\text{id} \times f)^*\mathcal{M}^{\otimes n})
& =
\chi(Y_t,
\mathcal{N}_t^{\otimes n} \otimes_{\mathcal{O}_{Y_t}}
f_t^*\mathcal{L}_t^{\otimes n}) \\
& =
n\deg(\mathcal{N}_t) + n\deg(f_t^*\mathcal{L}_t) +
\chi(Y_t, \mathcal{O}_{Y_t}) \\
& =
Q(n) + n\deg(\mathcal{L}_t) \\
& =
Q(n) + P(n) - P(0)
\end{align*}
这里使用了真曲线的 Riemann--Roch；更确切地说，使用了
《簇》中的定义 \ref{varieties-definition-degree-invertible-sheaf}、
引理 \ref{varieties-lemma-degree-tensor-product}，以及
$f_t$ 是同构这一事实。令 $P'(t) = Q(t) + P(t) - P(0)$，得到
$$
\mathit{Isom}_U(Y, X) =
\mathit{Isom}_U(Y, X) \cap \mathit{Mor}^{P', \mathcal{M}}_U(Y, X)
$$
该交集是 $\mathit{Mor}_U(Y, X)$ 的两个开子空间之交；见
《模叠》中的引理 \ref{moduli-lemma-Isom-in-Mor} 和注
\ref{moduli-remark-Mor-numerical}。
现在，$\mathit{Mor}^{P', \mathcal{M}}_U(Y, X)$
在 $U$ 上有限表示，故由《模叠》中的引理
\ref{moduli-lemma-Mor-qc-over-base}，它是 Noether 代数空间。
因此该交集也是 Noether 代数空间，证明完成。
\end{proof}

\begin{remark}
\label{moduli-curves-remark-boundedness-aut-does-not-work-surfaces}
引理 \ref{moduli-curves-lemma-curves-diagonal-separated-fp} 的证明中的有界性论证
对曲面的模空间并不适用；事实上，其结论也是错误的，例如域上的
K3 曲面可以具有无限离散自同构群。该论证不适用的“原因”是：
对于域上的射影曲面 $S$，给定 Hilbert 多项式分别为 $Q$ 和 $P$
的丰沛可逆层 $\mathcal{N}$ 和 $\mathcal{L}$，
$\mathcal{N} \otimes_{\mathcal{O}_S} \mathcal{L}$ 的 Hilbert 多项式
并不存在先验界。用相交论的语言来说，若 $H_1$、$H_2$ 是 $S$
上的丰沛有效 Cartier 除子，则相交数 $H_1 \cdot H_2$
不能由 $H_1 \cdot H_1$ 和 $H_2 \cdot H_2$ 给出（上）界。
\end{remark}

\begin{lemma}
\label{moduli-curves-lemma-curves-qs-lfp}
态射 $\Curvesstack \to \Spec(\mathbf{Z})$ 拟分离且局部有限表示。
\end{lemma}

\begin{proof}
要检验 $\Curvesstack \to \Spec(\mathbf{Z})$ 拟分离，需要证明其
对角态射拟紧且拟分离。这立即由引理
\ref{moduli-curves-lemma-curves-diagonal-separated-fp} 得出。
要证明 $\Curvesstack \to \Spec(\mathbf{Z})$ 局部有限表示，只需证明
$\Curvesstack$ 保持极限；见《叠的极限》中的命题
\ref{stacks-limits-proposition-characterize-locally-finite-presentation}。
这就是《Quot》中的引理 \ref{quot-lemma-curves-limits}。
\end{proof}





\section{曲线叠的开子叠}
\label{moduli-curves-section-open}

\noindent
下文常用代数空间态射的一个性质 $P$ 来刻画
$\Curvesstack$ 的开子叠。要说明 $P$ 定义开子叠，只需检验
\begin{enumerate}
\item[(o)] 给定曲线族 $f : X \to S$，存在最大开子概形
$S' \subset S$，使得
$f|_{f^{-1}(S')} : f^{-1}(S') \to S'$ 具有性质 $P$，
并且 $S'$ 的构造与任意基变换交换。
\end{enumerate}
事实上，假设 (o) 成立。取概形 $U$ 以及满光滑态射
$m : U \to \Curvesstack$。令
$R = U \times_{\Curvesstack} U$，并以 $t, s : R \to U$
表示投影。回忆 $\Curvesstack = [U/R]$ 是一个展示；见
《代数叠》中的引理 \ref{algebraic-lemma-stack-presentation} 和定义
\ref{algebraic-definition-presentation}。
由 $\Curvesstack$ 作为曲线叠的构造，态射 $m$ 是曲线族
$C \to U$ 的分类态射。图
$$
\xymatrix{
R \ar[r]_s \ar[d]_t & U \ar[d] \\
U \ar[r] & \Curvesstack
}
$$
的 $2$-交换性蕴含
$C \times_{U, s} R  \cong C \times_{U, t} R$
（这是 $R$ 上曲线族的同构）。令 $W \subset U$ 为满足如下条件的
最大开子概形：
% P12-ERR-0236: see ../control/ERRATA_CANDIDATES.jsonl
$f|_{f^{-1}(W)} : f^{-1}(W) \to W$ 具有 (o) 中的性质 $P$。
由于按 (o)，$W$ 的构造与基变换交换，再结合上述同构，得到
$s^{-1}(W) = t^{-1}(W)$。因此，由《叠的性质》中的引理
\ref{stacks-properties-lemma-immersion-presentation}，
$W \subset U$ 对应于开子叠
$$
\Curvesstack^P \subset \Curvesstack
$$

\medskip\noindent
% P12-ERR-0237: see ../control/ERRATA_CANDIDATES.jsonl
继续采用上一段的设定。我们断言开子叠 $\Curvesstack^P$
具有以下两个泛性质：
\begin{enumerate}
\item 给定曲线族 $X \to S$，下列条件等价：
\begin{enumerate}
\item 分类态射 $S \to \Curvesstack$ 经 $\Curvesstack^P$ 分解；
\item 态射 $X \to S$ 具有性质 $P$；
\end{enumerate}
% P12-ERR-0238: see ../control/ERRATA_CANDIDATES.jsonl
\item 给定域 $k$ 上维数 $\leq 1$ 的真概形 $X$，下列条件等价：
\begin{enumerate}
\item 分类态射 $\Spec(k) \to \Curvesstack$ 经
$\Curvesstack^P$ 分解；
\item 态射 $X \to \Spec(k)$ 具有性质 $P$。
\end{enumerate}
\end{enumerate}
这可通过考察 $2$-纤维积
$$
\xymatrix{
T \ar[r]_p \ar[d]_q & U \ar[d] \\
S \ar[r] & \Curvesstack
}
$$
得到。注意，$T \to S$ 是 $U \to \Curvesstack$ 的基变换，
故满且光滑。因此，(o) 给出的开集 $S' \subset S$
由其在 $T$ 中的逆像确定。然而，由 (o) 中这些开集在基变换下
的不变性，并且由 $2$-交换性有
$X \times_S T \cong C \times_U T$，得到 $T$ 中开集的等式
$q^{-1}(S') = p^{-1}(W)$。这立即蕴含 (1)，而 (2) 是 (1) 的特例。

\medskip\noindent
给定代数空间态射的两个性质 $P$ 和 $Q$，假设已经证明
$\Curvesstack^Q$ 是 $\Curvesstack$ 的开子叠，则可以用完全相同的
方法证明
$\Curvesstack^{Q, P} \subset \Curvesstack^Q$
是开子叠。我们略去精确说明。

\section{具有有限约化自同构群的曲线}
\label{moduli-curves-section-finite-aut}

\noindent
设 $X$ 是域 $k$ 上维数 $\leq 1$ 的真概形，即
$\Curvesstack$ 在 $k$ 上的对象。由引理
\ref{moduli-curves-lemma-curves-diagonal-separated-fp}，自同构群代数空间
$\mathit{Aut}(X)$ 在 $k$ 上有限型且分离。
特别地，$\mathit{Aut}(X)$ 是群概形；见《空间中群胚进阶》中的引理
\ref{spaces-more-groupoids-lemma-group-space-scheme-locally-finite-type-over-k}。
若 $k$ 的特征为零，则 $\mathit{Aut}(X)$ 约化，甚至在 $k$
上光滑（《群胚》中的引理
\ref{groupoids-lemma-group-scheme-characteristic-zero-smooth}）。
然而，一般而言，即使 $X$ 几何约化，$\mathit{Aut}(X)$
也不约化。

\begin{example}[非约化自同构群]
\label{moduli-curves-example-non-reduced}
设 $k$ 是特征 $2$ 的代数闭域。令
$Y = Z = \mathbf{P}^1_k$。在 $\mathbf{A}^1_k$ 中取三个两两不同的
$k$-值点 $a, b, c$。把
% P12-ERR-0239: see ../control/ERRATA_CANDIDATES.jsonl
$\mathbf{A}^1_k \subset \mathbf{P}^1_k = Y = Z$ 看作开子概形，
得到闭浸入
$$
T =  \Spec(k[t]/(t - a)^2) \amalg \Spec(k[t]/(t - b)^2)
\amalg \Spec(k[t]/(t - c)^2)
\longrightarrow
\mathbf{P}^1_k
$$
令 $X$ 为图
$$
\xymatrix{
T \ar[r] \ar[d] & Y \ar[d] \\
Z \ar[r] & X
}
$$
中的推出。令 $U \subset X$ 为
$\mathbf{A}^1_k \amalg \mathbf{A}^1_k$ 的像所构成的仿射开部分。
则有等化子图
$$
\xymatrix{
\mathcal{O}_X(U) \ar[r] &
k[t] \times k[t] \ar@<1ex>[r] \ar@<-1ex>[r] &
k[t]/(t - a)^2 \times k[t]/(t - b)^2 \times k[t]/(t - c)^2
}
$$
在双数 $A = k[\epsilon]$ 上，该等化子图有一个非平凡自同构，
它把 $t$ 送到 $t + \epsilon$。留给读者验证：该自同构扩张为
$X$ 在 $A$ 上的自同构。另一方面，读者容易证明 $X$ 在 $k$
上的自同构群有限。因此 $\mathit{Aut}(X)$ 必不约化。
\end{example}

\noindent
设 $X$ 是域 $k$ 上维数 $\leq 1$ 的真概形，即
$\Curvesstack$ 在 $k$ 上的对象。即使 $X$ 光滑且几何连通，
若 $\mathit{Aut}(X)$ 几何约化，也未必有维数 $0$。

\begin{example}[正维光滑自同构群]
\label{moduli-curves-example-pos-dim}
设 $k$ 是代数闭域。若 $X$ 是光滑的亏格 $0$、相应地\ $1$ 的曲线，
则其自同构群的维数为 $3$、相应地\ $1$。
事实上，在亏格 $0$ 的情形，由《代数曲线》中的命题
\ref{curves-proposition-projective-line}，有
$X \cong \mathbf{P}^1_k$。由于函子同构
$$
\mathit{Aut}(\mathbf{P}^1_k) = \text{PGL}_{2, k}
$$
可见维数为 $3$。另一方面，若 $X$ 的亏格为 $1$，则映射
$X = \underline{\Hilbfunctor}^1_{X/k} \to
\underline{\Picardfunctor}^1_{X/k}$ 是同构；见
《曲线的 Picard 概形》中的引理 \ref{pic-lemma-picard-pieces} 以及
《代数曲线》中的定理
\ref{curves-theorem-curves-rational-maps}。
因此 $X$ 具有 Abel 簇结构
（因为 $\underline{\Picardfunctor}^1_{X/k} \cong
\underline{\Picardfunctor}^0_{X/k}$）。
% P12-ERR-0240: see ../control/ERRATA_CANDIDATES.jsonl
特别地，$X$ 的（余）切丛平凡
（《群胚》中的引理
\ref{groupoids-lemma-group-scheme-module-differentials}）。
由此得到 $\dim_k H^0(X, T_X) = 1$，故
$\dim \mathit{Aut}(X) \leq 1$。另一方面，平移
（把 $X$ 看作群概形）给出 $\text{Aut}(X)$ 的一个 $1$-维部分，
% P12-ERR-0241: see ../control/ERRATA_CANDIDATES.jsonl
所以其维数确为 $1$。
\end{example}

\noindent
事实表明，$\Curvesstack$ 有一个开子叠，参数化自同构群
几何约化且有限的曲线。精确陈述如下。

\begin{lemma}
\label{moduli-curves-lemma-DM-curves}
% P12-ERR-0242: see ../control/ERRATA_CANDIDATES.jsonl
存在开子叠 $\Curvesstack^{DM} \subset \Curvesstack$，具有以下性质：
\begin{enumerate}
\item $\Curvesstack^{DM} \subset \Curvesstack$ 是最大的
DM 开子叠；
\item 给定曲线族 $X \to S$，下列条件等价：
\begin{enumerate}
\item 分类态射 $S \to \Curvesstack$ 经 $\Curvesstack^{DM}$ 分解；
\item 群代数空间 $\mathit{Aut}_S(X)$ 在 $S$ 上非分歧；
\end{enumerate}
% P12-ERR-0243: see ../control/ERRATA_CANDIDATES.jsonl
\item 给定域 $k$ 上维数 $\leq 1$ 的真概形 $X$，下列条件等价：
\begin{enumerate}
\item 分类态射 $\Spec(k) \to \Curvesstack$ 经
$\Curvesstack^{DM}$ 分解；
\item $\mathit{Aut}(X)$ 在 $k$ 上几何约化且维数为 $0$；
\item $\mathit{Aut}(X) \to \Spec(k)$ 非分歧。
\end{enumerate}
\end{enumerate}
\end{lemma}

\begin{proof}
具有性质 (1) 的开子叠之存在性由《叠的态射》中的引理
\ref{stacks-morphisms-lemma-open-DM-locus} 给出。
由《叠的态射》中的引理
\ref{stacks-morphisms-lemma-points-DM-locus}，该开子叠的点由
(3)(c) 刻画。(3)(b) 与 (3)(c) 的等价性是如下陈述：
若代数空间 $G$ 在域 $k$ 上局部有限型、几何约化且维数为 $0$，
则它在 $k$ 上非分歧。首先，由《域上的空间》中的引理
\ref{spaces-over-fields-lemma-locally-finite-type-dim-zero}，
$G$ 是概形。然后可在 $G$ 中取一个仿射开集，注意到它在 $k$
上为真，并应用《簇》中的引理
\ref{varieties-lemma-proper-geometrically-reduced-global-sections}。
略去次要细节。

\medskip\noindent
由 (3)，(2) 成立。事实上，态射
$\mathit{Aut}_S(X) \to S$ 局部有限型。因此，可以通过在概形 $S$
各点的纤维上检验，判断 $\mathit{Aut}_S(X) \to S$
是否在 $\mathit{Aut}_S(X)$ 的所有点处非分歧；见
《空间的态射》中的引理
\ref{spaces-morphisms-lemma-where-unramified}。
但基变换到 $S$ 的一点后，就回到 (3)(a) 与 (3)(c) 的等价性。
\end{proof}

\begin{lemma}
\label{moduli-curves-lemma-in-DM-locus-vector-fields}
设 $X$ 是域 $k$ 上维数 $\leq 1$ 的真概形。则性质
(3)(a)、(b)、(c) 还等价于
$\text{Der}_k(\mathcal{O}_X, \mathcal{O}_X) = 0$。
\end{lemma}

\begin{proof}
在上面的讨论中已经看到，$G = \mathit{Aut}(X)$ 是
$\Spec(k)$ 上有限型且分离的群概形；这里使用了引理
\ref{moduli-curves-lemma-curves-diagonal-separated-fp} 和《空间中群胚进阶》中的引理
\ref{spaces-more-groupoids-lemma-group-space-scheme-locally-finite-type-over-k}。
于是，$G$ 在 $k$ 上非分歧当且仅当 $\Omega_{G/k} = 0$
（《态射》中的引理 \ref{morphisms-lemma-unramified-omega-zero}）。
由《群胚》中的引理
\ref{groupoids-lemma-group-scheme-module-differentials}，
若 $T_{G/k, e} = 0$，则上述消失成立；这里 $T_{G/k, e}$ 是
$G$ 在单位元 $e \in G(k)$ 处的切空间，见《簇》中的定义
\ref{varieties-definition-tangent-space} 以及《簇》中的引理
\ref{varieties-lemma-tangent-space-cotangent-space} 里的公式。
由于 $\kappa(e) = k$，该切空间可用如下态射定义：
$\alpha : \Spec(k[\epsilon]) \to G = \mathit{Aut}(X)$，
其在 $\Spec(k)$ 上的限制为 $e$。于是，只需考察自同构
% P12-ERR-0244: see ../control/ERRATA_CANDIDATES.jsonl
$$
\alpha :
X \times_{\Spec(k)} \Spec(k[\epsilon])
\longrightarrow
X \times_{\Spec(k)} \Spec(k[\epsilon])
$$
它在 $\Spec(k[\epsilon])$ 上，且在 $\Spec(k)$ 上的限制为
$\text{id}_X$。这样的自同构称为无穷小自同构。

\medskip\noindent
$X$ 的无穷小自同构与 $\mathcal{O}_X$ 在 $k$ 上的导子按 $1$-对-$1$ 对应。
这由《态射进阶》中的引理
\ref{more-morphisms-lemma-difference-derivation} 和
\ref{more-morphisms-lemma-action-by-derivations} 得出
（我们只需要第一个引理，因为不关心反方向；另请参见
《态射进阶》中的注
\ref{more-morphisms-remark-another-special-case} 的说明）。
关于证明该等式的另一种论证，参见《形变问题》中的引理
\ref{examples-defos-lemma-schemes-TI}。
\end{proof}

\section{Cohen--Macaulay 曲线}
\label{moduli-curves-section-CM}

\noindent
$\Curvesstack$ 有一个开子叠，参数化 Cohen--Macaulay ``曲线''。

\begin{lemma}
\label{moduli-curves-lemma-CM-curves}
% P12-ERR-0245: see ../control/ERRATA_CANDIDATES.jsonl
存在开子叠 $\Curvesstack^{CM} \subset \Curvesstack$，使得：
\begin{enumerate}
\item 给定曲线族 $X \to S$，下列条件等价：
\begin{enumerate}
\item 分类态射 $S \to \Curvesstack$ 经 $\Curvesstack^{CM}$ 分解；
\item 态射 $X \to S$ 是 Cohen--Macaulay 态射；
\end{enumerate}
\item 给定域 $k$ 上的真概形 $X$，且 $\dim(X) \leq 1$，
下列条件等价：
\begin{enumerate}
\item 分类态射 $\Spec(k) \to \Curvesstack$ 经
$\Curvesstack^{CM}$ 分解；
\item $X$ 是 Cohen--Macaulay 概形。
\end{enumerate}
\end{enumerate}
\end{lemma}

\begin{proof}
设 $f : X \to S$ 是曲线族。由《空间的态射进阶》中的引理
\ref{spaces-more-morphisms-lemma-flat-finite-presentation-CM-open}，集合
$$
W = \{x \in |X| : f \text{ 在下述点为 Cohen--Macaulay：}x\}
$$
在 $|X|$ 中是开的，并且这个开集的构造与任意基变换交换。
由于 $f$ 为真，$S$ 的子集
$$
S' = S \setminus f(|X| \setminus W)
$$
是开的，且 $X \times_S S' \to S'$ 是 Cohen--Macaulay 态射。
此外，$S'$ 的构造与任意基变换交换，因为 $W$ 的构造具有此性质。
% P12-ERR-0246: see ../control/ERRATA_CANDIDATES.jsonl
因此，由第 \ref{moduli-curves-section-open} 节所讨论的方法，得到具有所需性质的开子叠。
\end{proof}

\begin{lemma}
\label{moduli-curves-lemma-CM-1-curves}
% P12-ERR-0247: see ../control/ERRATA_CANDIDATES.jsonl
存在开子叠 $\Curvesstack^{CM, 1} \subset \Curvesstack$，使得：
\begin{enumerate}
\item 给定曲线族 $X \to S$，下列条件等价：
\begin{enumerate}
\item 分类态射 $S \to \Curvesstack$ 经
$\Curvesstack^{CM, 1}$ 分解；
\item 态射 $X \to S$ 是 Cohen--Macaulay 态射，且具有相对维数 $1$
（《空间的态射》中的定义
\ref{spaces-morphisms-definition-relative-dimension}）；
\end{enumerate}
\item 给定域 $k$ 上的真概形 $X$，且 $\dim(X) \leq 1$，
下列条件等价：
\begin{enumerate}
\item 分类态射 $\Spec(k) \to \Curvesstack$ 经
$\Curvesstack^{CM, 1}$ 分解；
\item $X$ 是 Cohen--Macaulay 概形，且 $X$ 是维数为 $1$ 的等维概形。
\end{enumerate}
\end{enumerate}
\end{lemma}

\begin{proof}
由引理 \ref{moduli-curves-lemma-CM-curves}，若该开子叠存在，
则显然有 $\Curvesstack^{CM, 1} \subset \Curvesstack^{CM}$。
设 $f : X \to S$ 是曲线族，且 $f$ 是 Cohen--Macaulay 态射。
由《空间的态射进阶》中的引理
\ref{spaces-more-morphisms-lemma-lfp-CM-relative-dimension}，有按开闭子空间的分解
$$
X = X_0 \amalg X_1
$$
使得 $X_0 \to S$ 具有相对维数 $0$，而 $X_1 \to S$ 具有相对维数 $1$。
由于 $f$ 为真，$S$ 的子集
$$
S' = S \setminus f(|X_0|)
$$
是开的，且 $X \times_S S' \to S'$ 是相对维数为 $1$ 的
Cohen--Macaulay 态射。此外，$S'$ 的构造与任意基变换交换，
因为上述分解的构造具有此性质（相对维数在基变换下表现良好；见
《空间的态射》中的引理
\ref{spaces-morphisms-lemma-dimension-fibre-after-base-change}）。
因此，由第 \ref{moduli-curves-section-open} 节所讨论的方法，得到具有所需性质的开子叠。
\end{proof}

\section{给定亏格的曲线}
\label{moduli-curves-section-genus}

\noindent
Stacks Project 采用如下约定：域 $k$ 上的真 $1$-维概形 $X$
的亏格 $g$ 仅在 $H^0(X, \mathcal{O}_X) = k$ 时定义。
在此情况下，$g = \dim_k H^1(X, \mathcal{O}_X)$。
见《代数曲线》中的第 \ref{curves-section-genus} 节。
定义亏格所需的条件定义了一个开子叠；该开子叠又分解为开子叠的不交并，
每个亏格对应其中一个开子叠。

\begin{lemma}
\label{moduli-curves-lemma-pre-genus-curves}
% P12-ERR-0248: see ../control/ERRATA_CANDIDATES.jsonl
存在开子叠 $\Curvesstack^{h0, 1} \subset \Curvesstack$，使得：
\begin{enumerate}
\item 给定曲线族 $f : X \to S$，下列条件等价：
\begin{enumerate}
\item 分类态射 $S \to \Curvesstack$ 经
$\Curvesstack^{h0, 1}$ 分解；
\item $f_*\mathcal{O}_X = \mathcal{O}_S$，该等式在任意基变换后仍成立，
且 $f$ 的纤维维数为 $1$；
\end{enumerate}
\item 给定域 $k$ 上的真概形 $X$，且 $\dim(X) \leq 1$，
下列条件等价：
\begin{enumerate}
\item 分类态射 $\Spec(k) \to \Curvesstack$ 经
$\Curvesstack^{h0, 1}$ 分解；
\item $H^0(X, \mathcal{O}_X) = k$ 且 $\dim(X) = 1$。
\end{enumerate}
\end{enumerate}
\end{lemma}

\begin{proof}
给定曲线族 $X \to S$，由《空间的导出范畴》中的引理
\ref{spaces-perfect-lemma-jump-loci-geometric}，满足
$\kappa(s) = H^0(X_s, \mathcal{O}_{X_s})$ 的 $s \in S$
所成集合在 $S$ 中是开的。又由《空间的态射进阶》中的引理
\ref{spaces-more-morphisms-lemma-dimension-fibres-proper-flat}，
$S$ 中纤维维数为 $1$ 的点所成集合也是开的。
此外，若 $f : X \to S$ 是所有纤维维数均为 $1$ 的曲线族
（特别地，$f$ 是满射），则条件 (1)(b) 等价于对每个 $s \in S$ 都有
$\kappa(s) = H^0(X_s, \mathcal{O}_{X_s})$；见《空间的导出范畴》中的引理
\ref{spaces-perfect-lemma-proper-flat-h0}。
因此，该引理由第 \ref{moduli-curves-section-open} 节的一般讨论得出。
\end{proof}

\begin{lemma}
\label{moduli-curves-lemma-pre-genus-in-CM-1}
作为 $\Curvesstack$ 的开子叠，有
$\Curvesstack^{h0, 1} \subset \Curvesstack^{CM, 1}$。
\end{lemma}

\begin{proof}
见《代数曲线》中的引理 \ref{curves-lemma-automatic} 以及引理
\ref{moduli-curves-lemma-pre-genus-curves} 和 \ref{moduli-curves-lemma-CM-1-curves}。
\end{proof}

\begin{lemma}
\label{moduli-curves-lemma-genus}
设 $f : X \to S$ 是曲线族，并且对所有 $s \in S$ 都有
$\kappa(s) = H^0(X_s, \mathcal{O}_{X_s})$；换言之，分类态射
$S \to \Curvesstack$ 经 $\Curvesstack^{h0, 1}$ 分解
（引理 \ref{moduli-curves-lemma-pre-genus-curves}）。则：
\begin{enumerate}
\item $f_*\mathcal{O}_X = \mathcal{O}_S$，且该等式普遍成立；
\item $R^1f_*\mathcal{O}_X$ 是有限局部自由 $\mathcal{O}_S$-模；
\item 对任意态射 $h : S' \to S$，若 $f' : X' \to S'$ 是其基变换，
则 $h^*(R^1f_*\mathcal{O}_X) = R^1f'_*\mathcal{O}_{X'}$。
\end{enumerate}
\end{lemma}

\begin{proof}
应用《空间的导出范畴》中的引理
\ref{spaces-perfect-lemma-proper-flat-h0}。这证明了 (1)。
它还蕴含：在 $S$ 上局部可写成
$Rf_*\mathcal{O}_X = \mathcal{O}_S \oplus P$，其中 $P$
是 Tor 振幅位于 $[1, \infty)$ 的完美对象。回忆
$Rf_*\mathcal{O}_X$ 的构造与任意基变换交换
（《空间的导出范畴》中的引理
\ref{spaces-perfect-lemma-flat-proper-perfect-direct-image-general}）。
因此，对 $s \in S$，有
$$
H^i(P \otimes_{\mathcal{O}_S}^\mathbf{L} \kappa(s)) =
H^i(X_s, \mathcal{O}_{X_s})
\text{，其中 }i \geq 1
$$
由于 $X_s$ 是 $1$-维 Noether 概形，除非 $i = 1$，上式为零；见
《上同调》中的命题
\ref{cohomology-proposition-vanishing-Noetherian}。
于是 $P = H^1(P)[-1]$，并且 $H^1(P)$ 有限局部自由；例如可由
《代数进阶》中的引理
\ref{more-algebra-lemma-lift-perfect-from-residue-field} 得出。
由于所有构造都与基变换相容，(3) 也成立。
\end{proof}

\begin{lemma}
\label{moduli-curves-lemma-pre-genus-one-piece-per-genus}
存在开闭子叠的不交并分解
$$
\Curvesstack^{h0, 1} = \coprod\nolimits_{g \geq 0} \Curvesstack_g
$$
其中每个 $\Curvesstack_g$ 由下列条件刻画：
\begin{enumerate}
\item 给定曲线族 $f : X \to S$，下列条件等价：
\begin{enumerate}
\item 分类态射 $S \to \Curvesstack$ 经 $\Curvesstack_g$ 分解；
\item $f_*\mathcal{O}_X = \mathcal{O}_S$，该等式在任意基变换后仍成立，
$f$ 的纤维维数为 $1$，且 $R^1f_*\mathcal{O}_X$ 是秩为 $g$
的局部自由 $\mathcal{O}_S$-模；
\end{enumerate}
\item 给定域 $k$ 上的真概形 $X$，且 $\dim(X) \leq 1$，
下列条件等价：
\begin{enumerate}
\item 分类态射 $\Spec(k) \to \Curvesstack$ 经 $\Curvesstack_g$ 分解；
\item $\dim(X) = 1$、$k = H^0(X, \mathcal{O}_X)$，且 $X$
的亏格为 $g$。
\end{enumerate}
\end{enumerate}
\end{lemma}

\begin{proof}
由引理 \ref{moduli-curves-lemma-pre-genus-curves}，已经知道
$\Curvesstack^{h0, 1}$ 作为 $\Curvesstack$ 的开子叠存在，
并由本引理中不涉及 $R^1f_*$ 或 $H^1$ 的条件刻画。
由第 \ref{moduli-curves-section-open} 节的讨论和引理 \ref{moduli-curves-lemma-genus}，
立即得到开闭子叠分解的存在性。这证明了 (1) 中的刻画。
(2) 中的刻画由《代数曲线》中的定义
\ref{curves-definition-genus} 得出。
\end{proof}

\section{几何约化曲线}
\label{moduli-curves-section-geometrically-reduced}

\noindent
$\Curvesstack$ 有一个开子叠，参数化几何约化的 ``曲线''。

\begin{lemma}
\label{moduli-curves-lemma-geometrically-reduced-curves}
% P12-ERR-0249: see ../control/ERRATA_CANDIDATES.jsonl
存在开子叠 $\Curvesstack^{geomred} \subset \Curvesstack$，使得：
\begin{enumerate}
\item 给定曲线族 $X \to S$，下列条件等价：
\begin{enumerate}
\item 分类态射 $S \to \Curvesstack$ 经
$\Curvesstack^{geomred}$ 分解；
\item 态射 $X \to S$ 的纤维几何约化
（《空间的态射进阶》中的定义
\ref{spaces-more-morphisms-definition-geometrically-reduced-fibre}）；
\end{enumerate}
\item 给定域 $k$ 上的真概形 $X$，且 $\dim(X) \leq 1$，
下列条件等价：
\begin{enumerate}
\item 分类态射 $\Spec(k) \to \Curvesstack$ 经
$\Curvesstack^{geomred}$ 分解；
\item $X$ 在 $k$ 上几何约化。
\end{enumerate}
\end{enumerate}
\end{lemma}

\begin{proof}
设 $f : X \to S$ 是曲线族。由《空间的态射进阶》中的引理
\ref{spaces-more-morphisms-lemma-geometrically-reduced-open}，集合
$$
E = \{s \in S : \text{态射 }X \to S\text{ 在 }s
\text{ 处的纤维几何约化}\}
$$
在 $S$ 中是开的。由《空间的态射进阶》中的引理
\ref{spaces-more-morphisms-lemma-base-change-fibres-geometrically-reduced}，
这个开集的构造与任意基变换交换。因此，由第
\ref{moduli-curves-section-open} 节所讨论的方法，得到具有所需性质的开子叠。
\end{proof}

\begin{lemma}
\label{moduli-curves-lemma-geomred-in-CM}
作为 $\Curvesstack$ 的开子叠，有
$\Curvesstack^{geomred} \subset \Curvesstack^{CM}$。
\end{lemma}

\begin{proof}
这是因为维数 $\leq 1$ 的约化 Noether 概形是 Cohen--Macaulay 概形。
见《代数》中的引理 \ref{algebra-lemma-criterion-reduced}。
\end{proof}

\section{几何约化且连通的曲线}
\label{moduli-curves-section-geometrically-reduced-connected}

\noindent
$\Curvesstack$ 有一个开子叠，参数化几何约化且连通的 ``曲线''。
我们将立即排除 $0$-维对象。

\begin{lemma}
\label{moduli-curves-lemma-geometrically-reduced-connected-1-curves}
% P12-ERR-0250: see ../control/ERRATA_CANDIDATES.jsonl
存在开子叠 $\Curvesstack^{grc, 1} \subset \Curvesstack$，使得：
\begin{enumerate}
\item 给定曲线族 $X \to S$，下列条件等价：
\begin{enumerate}
\item 分类态射 $S \to \Curvesstack$ 经
$\Curvesstack^{grc, 1}$ 分解；
\item 态射 $X \to S$ 的几何纤维约化、连通且维数为 $1$；
\end{enumerate}
\item 给定域 $k$ 上的真概形 $X$，且 $\dim(X) \leq 1$，
下列条件等价：
\begin{enumerate}
\item 分类态射 $\Spec(k) \to \Curvesstack$ 经
$\Curvesstack^{grc, 1}$ 分解；
\item $X$ 几何约化、几何连通且维数为 $1$。
\end{enumerate}
\end{enumerate}
\end{lemma}

\begin{proof}
由引理 \ref{moduli-curves-lemma-geometrically-reduced-curves}、
\ref{moduli-curves-lemma-geomred-in-CM}、\ref{moduli-curves-lemma-CM-curves} 和
\ref{moduli-curves-lemma-CM-1-curves}，若所求开子叠存在，则显然有
$$
\Curvesstack^{grc, 1}
\subset
\Curvesstack^{geomred} \cap \Curvesstack^{CM, 1}
$$
设 $f : X \to S$ 是曲线族，其中 $f$ 是 Cohen--Macaulay 态射，
具有几何约化纤维且相对维数为 $1$。由《空间的态射进阶》中的引理
\ref{spaces-more-morphisms-lemma-stein-factorization-etale}，在 Stein 分解
$$
X \to T \to S
$$
中，态射 $T \to S$ 是 \'etale 态射。这蕴含存在开闭子概形
$S' \subset S$，使得 $X \times_S S' \to S'$ 具有几何连通纤维
（在《态射》中的引理 \ref{morphisms-lemma-finite-locally-free}
对有限局部自由态射 $T \to S$ 给出的分解中，这对应于 $S_1$）。
这个开集的构造与任意基变换交换，因为几何纤维的连通分支数在基变换下不变
（在当前特殊情形中，Stein 分解本身也与基变换交换，但得出结论并不需要这一点）。
因此，由第 \ref{moduli-curves-section-open} 节所讨论的方法，得到具有所需性质的开子叠。
\end{proof}

\begin{lemma}
\label{moduli-curves-lemma-geomredcon-in-h0-1}
作为 $\Curvesstack$ 的开子叠，有
$\Curvesstack^{grc, 1} \subset \Curvesstack^{h0, 1}$。
特别地，给定几何纤维约化、连通且维数为 $1$ 的曲线族
$f : X \to S$，则
% P12-ERR-0251: see ../control/ERRATA_CANDIDATES.jsonl
$R^1f_*\mathcal{O}_X$ 是有限局部自由 $\mathcal{O}_S$-模，
且其构造与任意基变换交换。
\end{lemma}

\begin{proof}
这由《簇》中的引理
\ref{varieties-lemma-proper-geometrically-reduced-global-sections} 以及引理
\ref{moduli-curves-lemma-pre-genus-curves} 和
\ref{moduli-curves-lemma-geometrically-reduced-connected-1-curves} 得出。
最后一个陈述由引理 \ref{moduli-curves-lemma-genus} 得出。
\end{proof}

\begin{lemma}
\label{moduli-curves-lemma-one-piece-per-genus}
存在开闭子叠的不交并分解
$$
\Curvesstack^{grc, 1} = \coprod\nolimits_{g \geq 0} \Curvesstack^{grc, 1}_g
$$
其中每个 $\Curvesstack^{grc, 1}_g$ 由下列条件刻画：
\begin{enumerate}
\item 给定曲线族 $f : X \to S$，下列条件等价：
\begin{enumerate}
\item 分类态射 $S \to \Curvesstack$ 经
$\Curvesstack^{grc, 1}_g$ 分解；
\item 态射 $f : X \to S$ 的几何纤维约化、连通且维数为 $1$，并且
$R^1f_*\mathcal{O}_X$ 是秩为 $g$ 的局部自由 $\mathcal{O}_S$-模；
\end{enumerate}
\item 给定域 $k$ 上的真概形 $X$，且 $\dim(X) \leq 1$，
下列条件等价：
\begin{enumerate}
\item 分类态射 $\Spec(k) \to \Curvesstack$ 经
$\Curvesstack^{grc, 1}_g$ 分解；
\item $X$ 几何约化、几何连通、维数为 $1$ 且亏格为 $g$。
\end{enumerate}
\end{enumerate}
\end{lemma}

\begin{proof}
第一种证明：令
$\Curvesstack^{grc, 1}_g = \Curvesstack^{grc, 1} \cap \Curvesstack_g$，
再结合引理 \ref{moduli-curves-lemma-geomredcon-in-h0-1} 和
\ref{moduli-curves-lemma-pre-genus-one-piece-per-genus}。
第二种证明：由第 \ref{moduli-curves-section-open} 节的讨论和引理
\ref{moduli-curves-lemma-geomredcon-in-h0-1}，立即得到开闭子叠分解的存在性。
这证明了 (1) 中的刻画。由于几何约化且连通的真 $1$-维概形 $X/k$
的亏格有定义（《代数曲线》中的定义 \ref{curves-definition-genus}
和《簇》中的引理
\ref{varieties-lemma-proper-geometrically-reduced-global-sections}），
且等于 $\dim_k H^1(X, \mathcal{O}_X)$，故 (2) 中的刻画也成立。
\end{proof}

\section{Gorenstein 曲线}
\label{moduli-curves-section-gorenstein}

\noindent
$\Curvesstack$ 有一个开子叠，参数化 Gorenstein ``曲线''。

\begin{lemma}
\label{moduli-curves-lemma-gorenstein-curves}
% P12-ERR-0252: see ../control/ERRATA_CANDIDATES.jsonl
存在开子叠 $\Curvesstack^{Gorenstein} \subset \Curvesstack$，使得：
\begin{enumerate}
\item 给定曲线族 $X \to S$，下列条件等价：
\begin{enumerate}
\item 分类态射 $S \to \Curvesstack$ 经
$\Curvesstack^{Gorenstein}$ 分解；
\item 态射 $X \to S$ 是 Gorenstein 态射；
\end{enumerate}
\item 给定域 $k$ 上的真概形 $X$，且 $\dim(X) \leq 1$，
下列条件等价：
\begin{enumerate}
\item 分类态射 $\Spec(k) \to \Curvesstack$ 经
$\Curvesstack^{Gorenstein}$ 分解；
\item $X$ 是 Gorenstein 概形。
\end{enumerate}
\end{enumerate}
\end{lemma}

\begin{proof}
设 $f : X \to S$ 是曲线族。由《空间的态射进阶》中的引理
\ref{spaces-more-morphisms-lemma-flat-finite-presentation-gorenstein-open}，集合
$$
W = \{x \in |X| : f \text{ 在下述点是 Gorenstein：}x\}
$$
在 $|X|$ 中是开的，并且这个开集的构造与任意基变换交换。
由于 $f$ 为真，$S$ 的子集
$$
S' = S \setminus f(|X| \setminus W)
$$
是开的，且 $X \times_S S' \to S'$ 是 Gorenstein 态射。
此外，$S'$ 的构造与任意基变换交换，因为 $W$ 的构造具有此性质。
% P12-ERR-0253: see ../control/ERRATA_CANDIDATES.jsonl
因此，由第 \ref{moduli-curves-section-open} 节所讨论的方法，得到具有所需性质的开子叠。
\end{proof}

\begin{lemma}
\label{moduli-curves-lemma-gorenstein-1-curves}
% P12-ERR-0254: see ../control/ERRATA_CANDIDATES.jsonl
存在开子叠
$\Curvesstack^{Gorenstein, 1} \subset \Curvesstack$，使得：
\begin{enumerate}
\item 给定曲线族 $X \to S$，下列条件等价：
\begin{enumerate}
\item 分类态射 $S \to \Curvesstack$ 经
$\Curvesstack^{Gorenstein, 1}$ 分解；
\item 态射 $X \to S$ 是 Gorenstein 态射，且具有相对维数 $1$
（《空间的态射》中的定义
\ref{spaces-morphisms-definition-relative-dimension}）；
\end{enumerate}
\item 给定域 $k$ 上的真概形 $X$，且 $\dim(X) \leq 1$，
下列条件等价：
\begin{enumerate}
\item 分类态射 $\Spec(k) \to \Curvesstack$ 经
$\Curvesstack^{Gorenstein, 1}$ 分解；
\item $X$ 是 Gorenstein 概形，且 $X$ 是维数为 $1$ 的等维概形。
\end{enumerate}
\end{enumerate}
\end{lemma}

\begin{proof}
回忆 Gorenstein 概形是 Cohen--Macaulay 概形
（《概形的对偶性》中的引理 \ref{duality-lemma-gorenstein-CM}），并且
% P12-ERR-0255: see ../control/ERRATA_CANDIDATES.jsonl
Gorenstein 态射是 Cohen--Macaulay 态射
（《概形的对偶性》中的引理 \ref{duality-lemma-gorenstein-CM-morphism}）。
因此，可以令 $\Curvesstack^{Gorenstein, 1}$ 等于
$\Curvesstack$ 内 $\Curvesstack^{Gorenstein}$ 与
$\Curvesstack^{CM, 1}$ 的交，并使用引理
\ref{moduli-curves-lemma-gorenstein-curves} 和 \ref{moduli-curves-lemma-CM-1-curves}。
\end{proof}

\section{局部完全交曲线}
\label{moduli-curves-section-lci}

\noindent
$\Curvesstack$ 有一个开子叠，参数化局部完全交 ``曲线''。

\begin{lemma}
\label{moduli-curves-lemma-lci-curves}
% P12-ERR-0256: see ../control/ERRATA_CANDIDATES.jsonl
存在开子叠 $\Curvesstack^{lci} \subset \Curvesstack$，使得：
\begin{enumerate}
\item 给定曲线族 $X \to S$，下列条件等价：
\begin{enumerate}
\item 分类态射 $S \to \Curvesstack$ 经 $\Curvesstack^{lci}$ 分解；
\item $X \to S$ 是局部完全交态射；
\item $X \to S$ 是 syntomic 态射。
\end{enumerate}
% P12-ERR-0257: see ../control/ERRATA_CANDIDATES.jsonl
\item 给定域 $k$ 上维数 $\leq 1$ 的真概形 $X$，下列条件等价：
\begin{enumerate}
\item 分类态射 $\Spec(k) \to \Curvesstack$ 经
$\Curvesstack^{lci}$ 分解；
\item $X$ 在 $k$ 上是局部完全交。
\end{enumerate}
\end{enumerate}
\end{lemma}

\begin{proof}
回忆 syntomic 态射等价于平坦的局部完全交态射；见
《空间的态射进阶》中的引理
\ref{spaces-more-morphisms-lemma-flat-lci}。因此 (1)(b) 等价于 (1)(c)。
在第 \ref{moduli-curves-section-open} 节中已经看到，只需证明：给定曲线族
$f : X \to S$，存在开子概形 $S' \subset S$，使得
$S' \times_S X \to S'$ 是局部完全交态射，并且 $S'$ 的构造
与任意基变换交换。这由更一般的《空间的态射进阶》中的引理
\ref{spaces-more-morphisms-lemma-where-lci} 得出。
\end{proof}

\section{具有孤立奇点的曲线}
\label{moduli-curves-section-curves-isolated}

\noindent
可以考察 $\Curvesstack$ 中参数化仅有有限多个奇点的 ``曲线''
所构成的开子叠（在我们的设定中，这些奇点可能对应于 $0$-维分支）。

\begin{lemma}
\label{moduli-curves-lemma-isolated-sings-curves}
% P12-ERR-0258: see ../control/ERRATA_CANDIDATES.jsonl
存在开子叠 $\Curvesstack^{+} \subset \Curvesstack$，使得：
\begin{enumerate}
\item 给定曲线族 $X \to S$，下列条件等价：
\begin{enumerate}
\item 分类态射 $S \to \Curvesstack$ 经 $\Curvesstack^{+}$ 分解；
\item 赋予 $X \to S$ 的奇异轨迹任意一个（或某一个）闭子空间结构后，
它在 $S$ 上有限。
\end{enumerate}
% P12-ERR-0259: see ../control/ERRATA_CANDIDATES.jsonl
\item 给定域 $k$ 上维数 $\leq 1$ 的真概形 $X$，下列条件等价：
\begin{enumerate}
\item 分类态射 $\Spec(k) \to \Curvesstack$ 经 $\Curvesstack^{+}$ 分解；
\item $X \to \Spec(k)$ 除有限多个点外处处光滑。
\end{enumerate}
\end{enumerate}
\end{lemma}

\begin{proof}
为证明该引理，只需证明：给定曲线族 $f : X \to S$，存在开子概形
$S' \subset S$，使得 $S' \times_S X \to S'$ 的纤维具有性质 (2)。
% P12-ERR-0260: see ../control/ERRATA_CANDIDATES.jsonl
（该开集的构造将自动与基变换交换。）按定义，$X \to S$
不光滑的点所成轨迹 $T \subset |X|$ 是闭的。令 $Z \subset X$
为在 $T$ 上赋予约化诱导代数空间结构所得到的闭子空间
（《空间的性质》中的定义
\ref{spaces-properties-definition-reduced-induced-space}）。
现在，若 $s \in S$ 是使 $Z_s$ 有限的点，则存在 $s$ 的开邻域
$U_s \subset S$，使得 $Z \cap f^{-1}(U_s) \to U_s$ 有限；见
《空间的态射进阶》中的引理
\ref{spaces-more-morphisms-lemma-proper-finite-fibre-finite-in-neighbourhood}。
这证明了该引理。
\end{proof}

\section{曲线叠的光滑轨迹}
\label{moduli-curves-section-smooth}

\noindent
态射
$$
\Curvesstack \longrightarrow \Spec(\mathbf{Z})
$$
在一个最大的开子叠上光滑；见《叠的态射》中的引理
\ref{stacks-morphisms-lemma-where-smooth}。我们希望给出曲线属于该轨迹的判据。
为此将使用少量形变理论。

\medskip\noindent
设 $k$ 是域，$X$ 是 $k$ 上维数 $\leq 1$ 的真概形。
为 $k$ 选取一个 Cohen 环 $\Lambda$；见《代数》中的引理
\ref{algebra-lemma-cohen-rings-exist}。于是处于《形变问题》中的例
\ref{examples-defos-example-schemes} 和引理
\ref{examples-defos-lemma-schemes-RS} 所描述的情形。因此，在剩余域为
$k$ 的 Artin 局部 $\Lambda$-代数所成范畴 $\mathcal{C}_\Lambda$ 上，
得到形变范畴 $\Deformationcategory_X$。

\begin{lemma}
\label{moduli-curves-lemma-in-smooth-locus}
在上述情形中，下列条件等价：
\begin{enumerate}
\item 分类态射 $\Spec(k) \to \Curvesstack$ 经
$\Curvesstack \to \Spec(\mathbf{Z})$ 光滑的开集分解；
\item 形变范畴 $\Deformationcategory_X$ 是无阻碍的。
\end{enumerate}
\end{lemma}

\begin{proof}
由于 $\Curvesstack \longrightarrow \Spec(\mathbf{Z})$ 局部有限表示
（引理 \ref{moduli-curves-lemma-curves-qs-lfp}），$\Curvesstack \longrightarrow
\Spec(\mathbf{Z})$ 光滑的开子叠之构造与平坦基变换交换
（《叠的态射》中的引理 \ref{stacks-morphisms-lemma-where-smooth}）。
由于 Cohen 环 $\Lambda$ 在 $\mathbf{Z}$ 上平坦，可以在 $\Lambda$ 上工作。
换言之，要证明
$$
\Lambda\text{-}\Curvesstack \longrightarrow \Spec(\Lambda)
$$
在由 $X/k$ 定义的点
$x_0 : \Spec(k) \to \Lambda\text{-}\Curvesstack$ 的一个开邻域上光滑，
当且仅当 $\Deformationcategory_X$ 无阻碍。

\medskip\noindent
该引理现在由《叠的几何》中的引理
\ref{stacks-geometry-lemma-characterize-smoothness} 和等式
$$
\Deformationcategory_X =
\mathcal{F}_{\Lambda\text{-}\Curvesstack, k, x_0}
$$
得出。建立这个等式并非完全平凡。事实上，左边是分类 $X$
作为 $A \in \Ob(\mathcal{C}_\Lambda)$ 上概形的所有平坦形变
$Y \to \Spec(A)$ 的形变范畴。右边是分类特殊纤维为 $X$
的所有平坦态射 $Y \to \Spec(A)$ 的形变范畴，其中 $Y$ 是代数空间，
且 $Y \to \Spec(A)$ 为真、有限表示并具有相对维数 $\leq 1$。
由于 $A$ 是 Artin 的，例如由《域上的空间》中的引理
\ref{spaces-over-fields-lemma-codim-1-point-in-schematic-locus} 可知
$Y$ 是概形。因此还需证明：$X$ 作为剩余域为 $k$ 的 Artin 局部环
$A$ 上概形的平坦形变 $Y \to \Spec(A)$ 为真、有限表示且相对维数
$\leq 1$。相对维数按纤维定义，故因为 $X/k$ 具有此性质，
$Y/A$ 也自动具有此性质。态射 $Y \to \Spec(A)$ 为真且局部有限表示，
因为 $X \to \Spec(k)$ 具有这些性质；见《态射进阶》中的引理
\ref{more-morphisms-lemma-deform-property}。
\end{proof}

\noindent
下面给出曲线叠中包含于光滑轨迹的一个 ``大'' 开子叠。

\begin{lemma}
\label{moduli-curves-lemma-big-smooth-part-curves}
开子叠
$$
\Curvesstack^{lci+} =
\Curvesstack^{lci} \cap \Curvesstack^{+}
\subset \Curvesstack
$$
具有以下性质：
\begin{enumerate}
\item $\Curvesstack^{lci+} \to \Spec(\mathbf{Z})$ 光滑；
\item 给定曲线族 $X \to S$，下列条件等价：
\begin{enumerate}
\item 分类态射 $S \to \Curvesstack$ 经 $\Curvesstack^{lci+}$ 分解；
\item $X \to S$ 是局部完全交态射，并且赋予 $X \to S$
的奇异轨迹任意一个（或某一个）闭子空间结构后，它在 $S$ 上有限；
\end{enumerate}
% P12-ERR-0261: see ../control/ERRATA_CANDIDATES.jsonl
\item 给定域 $k$ 上维数 $\leq 1$ 的真概形 $X$，下列条件等价：
\begin{enumerate}
\item 分类态射 $\Spec(k) \to \Curvesstack$ 经
$\Curvesstack^{lci+}$ 分解；
\item $X$ 在 $k$ 上是局部完全交，并且
$X \to \Spec(k)$ 除有限多个点外处处光滑。
\end{enumerate}
\end{enumerate}
\end{lemma}

\begin{proof}
若能证明存在开子叠 $\Curvesstack^{lci+}$，其点由 (2) 刻画，
则结合引理 \ref{moduli-curves-lemma-in-smooth-locus} 与《形变问题》中的引理
\ref{examples-defos-lemma-curve-isolated-lci}，可知 (1) 成立。由于
$$
\Curvesstack^{lci+} = \Curvesstack^{lci} \cap \Curvesstack^{+}
$$
在 $\Curvesstack$ 内成立，由引理 \ref{moduli-curves-lemma-lci-curves} 和
\ref{moduli-curves-lemma-isolated-sings-curves} 即得结论。
\end{proof}

\section{光滑曲线}
\label{moduli-curves-section-smooth-curves}

\noindent
本节研究 $\Curvesstack$ 中参数化光滑 ``曲线'' 的开子叠。

\begin{lemma}
\label{moduli-curves-lemma-smooth-curves}
% P12-ERR-0262: see ../control/ERRATA_CANDIDATES.jsonl
存在开子叠
$$
\Curvesstack^{smooth, 1} \subset \Curvesstack^{smooth} \subset \Curvesstack
$$
使得：
\begin{enumerate}
\item 给定曲线族 $f : X \to S$，下列条件等价：
\begin{enumerate}
\item 分类态射 $S \to \Curvesstack$ 经 $\Curvesstack^{smooth}$、
相应地，\ $\Curvesstack^{smooth, 1}$ 分解；
\item $f$ 光滑、相应地，\ 为相对维数 $1$ 的光滑态射；
\end{enumerate}
% P12-ERR-0263: see ../control/ERRATA_CANDIDATES.jsonl
\item 给定域 $k$ 上的真概形 $X$，且 $\dim(X) \leq 1$，
下列条件等价：
\begin{enumerate}
\item 分类态射 $\Spec(k) \to \Curvesstack$ 经 $\Curvesstack^{smooth}$、
相应地，\ $\Curvesstack^{smooth, 1}$ 分解；
\item $X$ 在 $k$ 上光滑；相应地，\ $X$ 在 $k$ 上光滑，且
$X$ 是维数为 $1$ 的等维概形。
\end{enumerate}
\end{enumerate}
\end{lemma}

\begin{proof}
为证明关于 $\Curvesstack^{smooth}$ 的陈述，只需证明：给定曲线族
$f : X \to S$，存在开子概形 $S' \subset S$，使得
$S' \times_S X \to S'$ 光滑，并且这个开集的构造与基变换交换。
已知存在最大的开集 $U \subset X$，使得 $U \to S$ 光滑，
而且 $U$ 的构造与任意基变换交换；见《空间的态射》中的引理
\ref{spaces-morphisms-lemma-where-smooth}。若
$T = |X| \setminus |U|$，则因为 $f$ 为真，$f(T)$ 在 $S$ 中闭。
令 $S' = S \setminus f(T)$，便得到所需开集。

\medskip\noindent
设 $f : X \to S$ 是曲线族，且 $f$ 光滑。则纤维 $X_s$
在 $\kappa(s)$ 上光滑，因而是 Cohen--Macaulay 概形
（例如可使用《代数》中的引理
\ref{algebra-lemma-smooth-over-field} 和
\ref{algebra-lemma-lci-CM} 看出）。因此，可以令
$$
\Curvesstack^{smooth, 1} = \Curvesstack^{smooth} \cap
\Curvesstack^{CM, 1}
$$
而所需等价性由已经对 $\Curvesstack^{smooth}$ 证明的结果以及引理
\ref{moduli-curves-lemma-CM-1-curves} 得出。
\end{proof}

\begin{lemma}
\label{moduli-curves-lemma-smooth-curves-smooth}
态射 $\Curvesstack^{smooth} \to \Spec(\mathbf{Z})$ 光滑。
\end{lemma}

\begin{proof}
由 $\Curvesstack^{smooth} \subset \Curvesstack^{lci+}$ 和引理
\ref{moduli-curves-lemma-big-smooth-part-curves} 立即得出。
\end{proof}

\begin{lemma}
\label{moduli-curves-lemma-smooth-curves-h0}
% P12-ERR-0264: see ../control/ERRATA_CANDIDATES.jsonl
存在开子叠 $\Curvesstack^{smooth, h0} \subset \Curvesstack$，使得：
\begin{enumerate}
\item 给定曲线族 $f : X \to S$，下列条件等价：
\begin{enumerate}
% P12-ERR-0265: see ../control/ERRATA_CANDIDATES.jsonl
\item 分类态射 $S \to \Curvesstack$ 经 $\Curvesstack^{smooth}$ 分解；
\item $f_*\mathcal{O}_X = \mathcal{O}_S$，该等式在任意基变换后仍成立，
且 $f$ 是相对维数为 $1$ 的光滑态射；
\end{enumerate}
% P12-ERR-0266: see ../control/ERRATA_CANDIDATES.jsonl
\item 给定域 $k$ 上的真概形 $X$，且 $\dim(X) \leq 1$，
下列条件等价：
\begin{enumerate}
\item 分类态射 $\Spec(k) \to \Curvesstack$ 经
$\Curvesstack^{smooth, h0}$ 分解；
\item $X$ 光滑，$\dim(X) = 1$，且 $k = H^0(X, \mathcal{O}_X)$；
\item $X$ 光滑，$\dim(X) = 1$，且 $X$ 几何连通；
\item $X$ 光滑，$\dim(X) = 1$，且 $X$ 几何整；
\item $X_{\overline{k}}$ 是光滑曲线。
\end{enumerate}
\end{enumerate}
\end{lemma}

\begin{proof}
若令
$$
\Curvesstack^{smooth, h0} = \Curvesstack^{smooth} \cap
\Curvesstack^{h0, 1}
$$
则由引理 \ref{moduli-curves-lemma-pre-genus-curves} 和
\ref{moduli-curves-lemma-smooth-curves} 可知 (1) 成立。事实上，这也给出
(2)(a) 与 (2)(b) 的等价性。为完成证明，需要说明 (2)(b)
分别等价于 (2)(c)、(2)(d) 和 (2)(e)。

\medskip\noindent
域上的光滑概形几何正规
（《簇》中的引理 \ref{varieties-lemma-smooth-geometrically-normal}），
光滑性在基变换下保持
（《态射》中的引理 \ref{morphisms-lemma-base-change-smooth}），
并且光滑性在目标上是 fpqc 局部的
（《下降》中的引理 \ref{descent-lemma-descending-property-smooth}）。
牢记这些事实，(2)(b)、(2)(c)、
% P12-ERR-0267: see ../control/ERRATA_CANDIDATES.jsonl
(2)(d) 和 (2)(e) 的等价性由《簇》中的引理
\ref{varieties-lemma-geometrically-normal-stein} 得出。
\end{proof}

\begin{definition}
\label{moduli-curves-definition-deligne-mumford-smooth}
\begin{reference}
\cite{DM}
\end{reference}
% P12-ERR-0269: see ../control/ERRATA_CANDIDATES.jsonl
我们用 $\mathcal{M}$ 表示引理 \ref{moduli-curves-lemma-smooth-curves-h0}
中引入的参数化曲线族的代数叠 $\Curvesstack^{smooth, h0}$，
并称其为{\it 光滑真曲线的模叠}。
对 $g \geq 0$，我们用 $\mathcal{M}_g$ 表示引理
\ref{moduli-curves-lemma-smooth-one-piece-per-genus} 中引入的代数叠，
并称其为{\it 亏格为 $g$ 的光滑真曲线的模叠}。
\end{definition}

\noindent
下面是必备引理。

\begin{lemma}
\label{moduli-curves-lemma-smooth-one-piece-per-genus}
存在开闭子叠的不交并分解
$$
\mathcal{M} = \coprod\nolimits_{g \geq 0} \mathcal{M}_g
$$
其中每个 $\mathcal{M}_g$ 由下列条件刻画：
\begin{enumerate}
\item 给定曲线族 $f : X \to S$，下列条件等价：
\begin{enumerate}
\item 分类态射 $S \to \Curvesstack$ 经 $\mathcal{M}_g$ 分解；
\item $X \to S$ 光滑，$f_*\mathcal{O}_X = \mathcal{O}_S$，
该等式在任意基变换后仍成立，且 $R^1f_*\mathcal{O}_X$
是秩为 $g$ 的局部自由 $\mathcal{O}_S$-模；
\end{enumerate}
% P12-ERR-0268: see ../control/ERRATA_CANDIDATES.jsonl
\item 给定域 $k$ 上的真概形 $X$，且 $\dim(X) \leq 1$，
下列条件等价：
\begin{enumerate}
\item 分类态射 $\Spec(k) \to \Curvesstack$ 经 $\mathcal{M}_g$ 分解；
\item $X$ 光滑，$\dim(X) = 1$，$k = H^0(X, \mathcal{O}_X)$，
且 $X$ 的亏格为 $g$；
\item $X$ 光滑，$\dim(X) = 1$，$X$ 几何连通，且 $X$ 的亏格为 $g$；
\item $X$ 光滑，$\dim(X) = 1$，$X$ 几何整，且 $X$ 的亏格为 $g$；
\item $X_{\overline{k}}$ 是亏格为 $g$ 的光滑曲线。
\end{enumerate}
\end{enumerate}
\end{lemma}

\begin{proof}
结合引理 \ref{moduli-curves-lemma-smooth-curves-h0} 和
\ref{moduli-curves-lemma-pre-genus-one-piece-per-genus}。也可以改用引理
\ref{moduli-curves-lemma-one-piece-per-genus}。
\end{proof}

\begin{lemma}
\label{moduli-curves-lemma-smooth-curves-h0-smooth}
态射 $\mathcal{M} \to \Spec(\mathbf{Z})$ 和
$\mathcal{M}_g \to \Spec(\mathbf{Z})$ 都光滑。
\end{lemma}

\begin{proof}
由于 $\mathcal{M}$ 是 $\Curvesstack^{lci+}$ 的开子叠，
该结论由引理 \ref{moduli-curves-lemma-big-smooth-part-curves} 得出。
\end{proof}

\section{光滑曲线的稠密性}
\label{moduli-curves-section-smooth-is-dense}

\noindent
本节标题有些误导，因为我们并不声称 $\Curvesstack^{smooth}$
在 $\Curvesstack$ 中稠密。事实上，这是错误的，Mumford 已在
\cite{PathologiesIV} 中证明这一点。不过，我们将看到光滑 ``曲线''
在一个大开集中稠密。

\begin{lemma}
\label{moduli-curves-lemma-smooth-dense}
包含关系
$$
|\Curvesstack^{smooth}| \subset |\Curvesstack^{lci+}|
$$
给出一个开稠密子集。
\end{lemma}

\begin{proof}
由《叠的性质》中的第 \ref{stacks-properties-section-points} 节所构造的
$|\Curvesstack^{lci+}|$ 上的拓扑，$|\Curvesstack^{smooth}|$
是开子集。设 $\xi \in |\Curvesstack^{lci+}|$ 是一点。
则存在域 $k$ 和 $k$ 上的概形 $X$，使得 $X$ 在 $k$ 上为真，
$\dim(X) \leq 1$，$X$ 在 $k$ 上是局部完全交，并且
% P12-ERR-0270: see ../control/ERRATA_CANDIDATES.jsonl
$X$ 在 $k$ 上除有限多个点外处处光滑，而且 $\xi$ 是由 $X$
确定的分类态射 $\Spec(k) \to \Curvesstack^{lci+}$ 的等价类。
见引理 \ref{moduli-curves-lemma-big-smooth-part-curves}。由《形变问题》中的引理
\ref{examples-defos-lemma-smoothing-proper-curve-isolated-lci}，
存在平坦射影态射 $Y \to \Spec(k[[t]])$，其泛纤维光滑，
其特殊纤维同构于 $X$。考虑由 $Y$ 确定的分类态射
$$
\Spec(k[[t]]) \longrightarrow \Curvesstack^{lci+}
$$
闭点的像为 $\xi$，而泛点的像属于 $|\Curvesstack^{smooth}|$。
由于泛点在 $|\Spec(k[[t]])|$ 中特化到闭点，故 $\xi$
属于 $|\Curvesstack^{smooth}|$ 的闭包，如所求。
\end{proof}

\section{结点曲线}
\label{moduli-curves-section-nodal-curves}

\noindent
结点曲线在代数几何中扮演特殊角色。建议读者简要浏览《代数曲线》中的
第 \ref{curves-section-nodal} 节和第
\ref{curves-section-families-nodal} 节，以及《空间的态射进阶》中的
第 \ref{spaces-more-morphisms-section-families-nodal} 节的部分讨论。

\begin{lemma}
\label{moduli-curves-lemma-nodal-curves}
% P12-ERR-0271: see ../control/ERRATA_CANDIDATES.jsonl
存在开子叠 $\Curvesstack^{nodal} \subset \Curvesstack$，使得：
\begin{enumerate}
\item 给定曲线族 $f : X \to S$，下列条件等价：
\begin{enumerate}
\item 分类态射 $S \to \Curvesstack$ 经 $\Curvesstack^{nodal}$ 分解；
\item $f$ 是相对维数为 $1$、至多有结点奇性的态射；
\end{enumerate}
% P12-ERR-0272: see ../control/ERRATA_CANDIDATES.jsonl
\item 给定域 $k$ 上的真概形 $X$，且 $\dim(X) \leq 1$，
下列条件等价：
\begin{enumerate}
\item 分类态射 $\Spec(k) \to \Curvesstack$ 经
$\Curvesstack^{nodal}$ 分解；
\item $X$ 的奇点至多为结点，且 $X$ 是维数为 $1$ 的等维概形。
\end{enumerate}
\end{enumerate}
\end{lemma}

\begin{proof}
事实上，只需证明：给定曲线族 $f : X \to S$，存在开子概形
$S' \subset S$，使得 $S' \times_S X \to S'$ 是相对维数为 $1$、
至多有结点奇性的态射，并且 $S'$ 的构造与任意基变换交换。
由《空间的态射进阶》中的引理
\ref{spaces-more-morphisms-lemma-locus-where-nodal}，存在最大的开子空间
$X' \subset X$，使得 $f|_{X'} : X' \to S$ 是相对维数为 $1$、
至多有结点奇性的态射。此外，$X'$ 的构造与基变换交换。
因此，可以取
$$
S' = S \setminus |f|(|X| \setminus |X'|)
$$
它是开的，因为真态射按定义普遍闭。
\end{proof}

\begin{lemma}
\label{moduli-curves-lemma-nodal-curves-smooth}
态射 $\Curvesstack^{nodal} \to \Spec(\mathbf{Z})$ 光滑。
\end{lemma}

\begin{proof}
由 $\Curvesstack^{nodal} \subset \Curvesstack^{lci+}$ 和引理
\ref{moduli-curves-lemma-big-smooth-part-curves} 立即得出。
\end{proof}

\section{相对对偶层}
\label{moduli-curves-section-relative-dualizing}

\noindent
本节主要用于在曲线族的情形引入记号。大部分工作已经在对偶性一章中完成。

\medskip\noindent
设 $f : X \to S$ 是曲线族。$D_\QCoh(\mathcal{O}_X)$ 中存在对象
$\omega_{X/S}^\bullet$，称为{\it 相对对偶复形}，具有如下性质：
对每个基变换图
$$
\xymatrix{
X_U \ar[d]_{f'} \ar[r]_{g'} & X \ar[d]^f \\
U \ar[r]^g & S
}
$$
若 $U = \Spec(A)$ 仿射，则复形
$\omega_{X_U/U}^\bullet = L(g')^*\omega_{X/S}^\bullet$ 表示函子
$$
D_\QCoh(\mathcal{O}_{X_U}) \longrightarrow \text{Mod}_A,\quad
K \longmapsto \Hom_U(Rf_*K, \mathcal{O}_U)
$$
更精确地，令 $(\omega_{X/S}^\bullet, \tau)$ 为《空间的对偶性》中的定义
\ref{spaces-duality-definition-relative-dualizing-proper-flat}
所定义的该曲线族的相对对偶复形。其存在性由《空间的对偶性》中的引理
\ref{spaces-duality-lemma-existence-relative-dualizing} 证明。
此外，$(\omega_{X/S}^\bullet, \tau)$ 的构造与任意基变换交换
（这本质上源于定义；精确参考是《空间的对偶性》中的引理
\ref{spaces-duality-lemma-base-change-relative-dualizing}）。
从现在起，我们将 $\omega_{X/S}^\bullet$ 的基变换与基变换后曲线族的
相对对偶复形等同，不再另行说明。

\medskip\noindent
设 $\{S_i \to S\}$ 是一个 \'etale 覆盖，其中 $S_i$ 仿射，且
$X_i = X \times_S S_i$ 是概形；见引理
\ref{moduli-curves-lemma-etale-locally-scheme}。由《空间的对偶性》中的引理
\ref{spaces-duality-lemma-compare}，$\omega_{X_i/S_i}^\bullet$
与《概形的对偶性》中的注
\ref{duality-remark-relative-dualizing-complex} 所讨论的概形间真、平坦、
有限表示态射 $f_i : X_i \to S_i$ 的相对对偶复形一致。
因此，为证明 $\omega_{X/S}^\bullet$ 的某个 \'etale 局部性质，
可以假设 $X \to S$ 是概形间态射，并使用概形对偶性一章中发展的理论。
更一般地，对 $X$ 的任意为概形的基变换，相对对偶复形都与
《概形的对偶性》中的注
\ref{duality-remark-relative-dualizing-complex} 的相对对偶复形一致。
从现在起，我们将使用这个等同，不再另行说明。

\medskip\noindent
特别地，设 $\Spec(k) \to S$ 是态射，其中 $k$ 是域。
% P12-ERR-0273: see ../control/ERRATA_CANDIDATES.jsonl
用 $X_k$ 表示该基变换（由《域上的空间》中的引理
\ref{spaces-over-fields-lemma-codim-1-point-in-schematic-locus}，
这是概形）。则 $\omega_{X_k/k}^\bullet$ 同构于《代数曲线》中的引理
\ref{curves-lemma-duality-dim-1} 的复形 $\omega_{X_k}^\bullet$
（两者表示同一个函子，故可使用 Yoneda 引理；而这实际上由上述评注得出）。
由此，余调层 $H^i(\omega_{X_k/k}^\bullet)$ 仅在 $i = 0, -1$
时可能非零。若 $X_k$ 是维数为 $1$ 的等维 Cohen--Macaulay 概形，
则仅有 $H^{-1}$；若 $X_k$ 还为 Gorenstein 概形，则
% P12-ERR-0274: see ../control/ERRATA_CANDIDATES.jsonl
$H^{-1}(\omega_{X_k/k})$ 可逆；见《代数曲线》中的引理
\ref{curves-lemma-duality-dim-1-CM} 和 \ref{curves-lemma-rr}。

\begin{lemma}
\label{moduli-curves-lemma-CM-dualizing}
设 $X \to S$ 是纤维为维数 $1$ 的等维 Cohen--Macaulay 概形的曲线族
（引理 \ref{moduli-curves-lemma-CM-1-curves}）。则
$\omega_{X/S}^\bullet = \omega_{X/S}[1]$，其中 $\omega_{X/S}$
是一个在 $S$ 上平坦的伪凝聚 $\mathcal{O}_X$-模，且其构造与任意基变换交换。
\end{lemma}

\begin{proof}
我们建议读者由上面对基变换到域后情形的讨论直接推出此结论。
我们的证明将使用一种稍显繁琐的归约，将问题化到 Noether 概形的情形。

\medskip\noindent
一旦证明 $\omega_{X/S}^\bullet = \omega_{X/S}[1]$ 且
$\omega_{X/S}$ 在 $S$ 上平坦，关于基变换的陈述就由
$\omega_{X/S}^\bullet$ 的构造与任意基变换交换这一已知事实得出。
此外，伪凝聚性将自动成立，因为按定义 $\omega_{X/S}^\bullet$
是伪凝聚的。其他余调层的消失与平坦性可在 \'etale 局部检验。
因此，可以假设 $f : X \to S$ 是概形间态射，且 $S$ 仿射
（见上面的讨论）。将 $S = \lim S_i$ 写成在 $\mathbf{Z}$ 上有限型的
仿射概形 $S_i$ 的余滤极限。由于 $\Curvesstack^{CM, 1}$ 在
$\mathbf{Z}$ 上局部有限表示（它是 $\Curvesstack$ 的开子叠；见引理
\ref{moduli-curves-lemma-CM-1-curves} 和 \ref{moduli-curves-lemma-curves-qs-lfp}），
可以找到一个 $i$ 以及曲线族 $X_i \to S_i$，其拉回为 $X \to S$
（《叠的极限》中的引理
\ref{stacks-limits-lemma-representable-by-spaces-limit-preserving}）。
必要时增大 $i$，可假设 $X_i$ 是概形；见《空间的极限》中的引理
\ref{spaces-limits-lemma-limit-is-scheme}。由于
$\omega_{X/S}^\bullet$ 的构造与任意基变换交换，可以用 $S_i$ 代替 $S$。
这样便可以并且确实假设 $S_i$ 是 Noether 概形。此时，由我们对纤维的假设，
$f$ 显然是 Cohen--Macaulay 态射
（《态射进阶》中的定义 \ref{more-morphisms-definition-CM}）。
此外，由《概形的对偶性》中的第
\ref{duality-section-upper-shriek} 节对 $f^!$ 的构造，有
$\omega_{X/S}^\bullet = f^!\mathcal{O}_S$。
% P12-ERR-0275: see ../control/ERRATA_CANDIDATES.jsonl
因此，该引理由《概形的对偶性》中的引理
\ref{duality-lemma-affine-flat-Noetherian-CM} 得出。
\end{proof}

\begin{definition}
\label{moduli-curves-definition-relative-dualizing-sheaf}
设 $f : X \to S$ 是纤维为维数 $1$ 的等维 Cohen--Macaulay 概形的曲线族
（引理 \ref{moduli-curves-lemma-CM-1-curves}）。则引理
\ref{moduli-curves-lemma-CM-dualizing} 中研究的 $\mathcal{O}_X$-模
$$
\omega_{X/S} = H^{-1}(\omega_{X/S}^\bullet)
$$
称为 $f$ 的{\it 相对对偶层}。
\end{definition}

\noindent
在定义 \ref{moduli-curves-definition-relative-dualizing-sheaf} 的情形，
相对对偶层 $\omega_{X/S}$ 具有以下性质（而且该性质在 $S$ 上局部刻画它）：
对每个基变换图
$$
\xymatrix{
X_U \ar[d]_{f'} \ar[r]_{g'} & X \ar[d]^f \\
U \ar[r]^g & S
}
$$
若 $U = \Spec(A)$ 仿射，则模
$\omega_{X_U/U} = (g')^*\omega_{X/S}$ 表示函子
$$
\QCoh(\mathcal{O}_{X_U}) \longrightarrow \text{Mod}_A,\quad
% P12-ERR-0276: see ../control/ERRATA_CANDIDATES.jsonl
\mathcal{F} \longmapsto \Hom_A(H^1(X, \mathcal{F}), A)
$$
这由上面给出的相对对偶复形的相应性质立即得出。特别地，若 $A = k$
是域，则恢复《代数曲线》中的引理
\ref{curves-lemma-duality-dim-1}、
\ref{curves-lemma-duality-dim-1-CM} 和 \ref{curves-lemma-rr}
所引入并研究的 $X_k$ 的对偶模。

\begin{lemma}
\label{moduli-curves-lemma-gorenstein-dualizing}
设 $X \to S$ 是纤维为维数 $1$ 的等维 Gorenstein 概形的曲线族
（引理 \ref{moduli-curves-lemma-gorenstein-1-curves}）。则相对对偶层
$\omega_{X/S}$ 是可逆 $\mathcal{O}_X$-模，且其构造与任意基变换交换。
\end{lemma}

\begin{proof}
这是因为由上面的讨论，相对对偶模到任一纤维的拉回均可逆。
或者，也可以完全仿照引理 \ref{moduli-curves-lemma-CM-dualizing} 的证明，
再由《概形的对偶性》中的引理
\ref{duality-lemma-affine-flat-Noetherian-gorenstein} 得到结论。
\end{proof}

\section{预稳定曲线}
\label{moduli-curves-section-prestable-curves}

\noindent
下述定义等价于通常公认的预稳定曲线族概念。

\begin{definition}
\label{moduli-curves-definition-prestable}
设 $f : X \to S$ 是曲线族。若满足下列条件，便称 $f$ 是一个
{\it 预稳定曲线族}：
\begin{enumerate}
\item $f$ 是相对维数为 $1$、至多有结点奇性的态射；并且
\item $f_*\mathcal{O}_X = \mathcal{O}_S$，而且该等式在任意基变换后仍成立
\footnote{事实上，只要求 $f_*\mathcal{O}_X = \mathcal{O}_S$ 即可，
因为此时 $f$ 的 Stein 分解是 \'etale 的；见《空间的态射进阶》中的引理
\ref{spaces-more-morphisms-lemma-stein-factorization-etale}。
也可以把该条件替换为要求几何纤维连通；见引理
\ref{moduli-curves-lemma-geomredcon-in-h0-1}。}。
\end{enumerate}
\end{definition}

\noindent
设 $X$ 是域 $k$ 上的真概形，且 $\dim(X) \leq 1$。
则 $X \to \Spec(k)$ 是曲线族，因此可以问它在上述定义的意义下是否预稳定
\footnote{此处不能使用“预稳定曲线”这一术语，因为“曲线”蕴含不可约性。
见《代数曲线》中的第 \ref{curves-section-families-nodal} 节的讨论。}。
展开定义可见，下列条件等价：
\begin{enumerate}
\item $X$ 预稳定；
\item $X$ 的奇点至多为结点，$\dim(X) = 1$，且
$k = H^0(X, \mathcal{O}_X)$；
\item $X_{\overline{k}}$ 连通，并且除有限多个结点外，
它在 $\overline{k}$ 上光滑
（《代数曲线》中的定义 \ref{curves-definition-multicross}）。
\end{enumerate}
这说明我们的定义与文献中的大多数定义一致。

\begin{lemma}
\label{moduli-curves-lemma-prestable-curves}
% P12-ERR-0277: see ../control/ERRATA_CANDIDATES.jsonl
存在开子叠 $\Curvesstack^{prestable} \subset \Curvesstack$，使得：
\begin{enumerate}
\item 给定曲线族 $f : X \to S$，下列条件等价：
\begin{enumerate}
\item 分类态射 $S \to \Curvesstack$ 经 $\Curvesstack^{prestable}$ 分解；
\item $X \to S$ 是预稳定曲线族；
\end{enumerate}
% P12-ERR-0278: see ../control/ERRATA_CANDIDATES.jsonl
\item 给定域 $k$ 上的真概形 $X$，且 $\dim(X) \leq 1$，
下列条件等价：
\begin{enumerate}
\item 分类态射 $\Spec(k) \to \Curvesstack$ 经
$\Curvesstack^{prestable}$ 分解；
\item $X$ 的奇点至多为结点，$\dim(X) = 1$，且
$k = H^0(X, \mathcal{O}_X)$。
\end{enumerate}
\end{enumerate}
\end{lemma}

\begin{proof}
给定曲线族 $X \to S$，它预稳定，当且仅当分类态射同时经
$\Curvesstack^{nodal}$ 和 $\Curvesstack^{h0, 1}$ 分解。
也可以使用 $\Curvesstack^{grc, 1}$（因为结点曲线几何约化，
所以它的 $H^0$ 等于基域，当且仅当它连通）。用公式表示即
$$
\Curvesstack^{prestable} =
\Curvesstack^{nodal} \cap \Curvesstack^{h0, 1} =
\Curvesstack^{nodal} \cap \Curvesstack^{grc, 1}
$$
因此，该引理由引理 \ref{moduli-curves-lemma-pre-genus-curves} 和
\ref{moduli-curves-lemma-nodal-curves} 得出。
\end{proof}

\noindent
对每个亏格 $g \geq 0$，都有分类亏格为 $g$ 的预稳定曲线的代数叠。
事实上，从现在起，我们称 $X \to S$ 是一个{\it 亏格为 $g$
的预稳定曲线族}，当且仅当分类态射 $S \to \Curvesstack$
经引理 \ref{moduli-curves-lemma-prestable-one-piece-per-genus} 的开子叠
$\Curvesstack^{prestable}_g$ 分解。

\begin{lemma}
\label{moduli-curves-lemma-prestable-one-piece-per-genus}
存在开闭子叠的分解
$$
\Curvesstack^{prestable} = \coprod\nolimits_{g \geq 0}
\Curvesstack^{prestable}_g
$$
其中每个 $\Curvesstack^{prestable}_g$ 由下列性质刻画：
\begin{enumerate}
\item 给定曲线族 $f : X \to S$，下列条件等价：
\begin{enumerate}
\item 分类态射 $S \to \Curvesstack$ 经
$\Curvesstack^{prestable}_g$ 分解；
\item $X \to S$ 是预稳定曲线族，且 $R^1f_*\mathcal{O}_X$
是秩为 $g$ 的局部自由 $\mathcal{O}_S$-模；
\end{enumerate}
\item 给定域 $k$ 上的真概形 $X$，且 $\dim(X) \leq 1$，
下列条件等价：
\begin{enumerate}
\item 分类态射 $\Spec(k) \to \Curvesstack$ 经
$\Curvesstack^{prestable}_g$ 分解；
\item $X$ 的奇点至多为结点，$\dim(X) = 1$，
$k = H^0(X, \mathcal{O}_X)$，且 $X$ 的亏格为 $g$。
\end{enumerate}
\end{enumerate}
\end{lemma}

\begin{proof}
我们已经看到 $\Curvesstack^{prestable}$ 包含于
$\Curvesstack^{h0, 1}$，故结论由引理
\ref{moduli-curves-lemma-prestable-curves} 和
\ref{moduli-curves-lemma-pre-genus-one-piece-per-genus} 得出。
\end{proof}

\begin{lemma}
\label{moduli-curves-lemma-prestable-curves-smooth}
态射 $\Curvesstack^{prestable} \to \Spec(\mathbf{Z})$ 和
$\Curvesstack^{prestable}_g \to \Spec(\mathbf{Z})$ 光滑。
\end{lemma}

\begin{proof}
由于 $\Curvesstack^{prestable}$ 是 $\Curvesstack^{nodal}$ 的开子叠，
结论由引理 \ref{moduli-curves-lemma-nodal-curves-smooth} 得出。
\end{proof}

\section{半稳定曲线}
\label{moduli-curves-section-semistable-curves}

\noindent
下述引理将帮助我们理解半稳定曲线族。

\begin{lemma}
\label{moduli-curves-lemma-semistable}
设 $f : X \to S$ 是亏格为 $g \geq 1$ 的预稳定曲线族。
设 $s \in S$ 是基概形上的一点，并设 $m \geq 2$。下列条件等价：
\begin{enumerate}
\item $X_s$ 没有有理尾
（《代数曲线》中的例 \ref{curves-example-rational-tail}）；
\item 对某个开子集 $s \in U \subset S$，态射
% P12-ERR-0279: see ../control/ERRATA_CANDIDATES.jsonl
$f^*f_*\omega_{X/S}^{\otimes m} \to \omega_{X/S}^{\otimes m}$
在 $f^{-1}(U)$ 上满射。
\end{enumerate}
\end{lemma}

\begin{proof}
假设 (2)。利用第 \ref{moduli-curves-section-relative-dualizing} 节的内容，
可知 $\omega_{X_s}^{\otimes m}$ 整体生成。然而，若
$C \subset X_s$ 是有理尾，则由《代数曲线》中的引理
\ref{curves-lemma-rational-tail-negative}，有
$\deg(\omega_{X_s}|_C) < 0$，故由《簇》中的引理
\ref{varieties-lemma-check-invertible-sheaf-trivial}，有
$H^0(C, \omega_{X_s}|_C) = 0$；这与它整体生成矛盾。
这证明了 (1)。

\medskip\noindent
假设 (1)。先假设 $g \geq 2$。由《代数曲线》中的引理
\ref{curves-lemma-contracting-rational-tails}，假设 (1) 蕴含
$\omega_{X_s}^{\otimes m}$ 整体生成。此外，由对偶性，有
$$
\Hom_{\kappa(s)}(H^1(X_s, \omega_{X_s}^{\otimes m}), \kappa(s)) =
H^0(X_s, \omega_{X_s}^{\otimes 1 - m})
$$
见《代数曲线》中的引理 \ref{curves-lemma-duality-dim-1-CM}。
由于 $\omega_{X_s}^{\otimes m}$ 整体生成，它在每个不可约分支上的限制
% P12-ERR-0280: see ../control/ERRATA_CANDIDATES.jsonl
次数非负。因此，$\omega_{X_s}^{\otimes 1 - m}$ 在每个不可约分支上的
限制次数非正。又由 Riemann--Roch（《代数曲线》中的引理
\ref{curves-lemma-rr}），有
$\deg(\omega_{X_s}^{\otimes 1 - m}) = (1 - m)(2g - 2) < 0$，
故由《簇》中的引理 \ref{varieties-lemma-no-sections-dual-nef}，
可知该 $H^0$ 为零。由上同调与基变换可知
$$
E = Rf_*\omega_{X/S}^{\otimes m}
$$
是完美复形，且其构造与任意基变换交换
（《空间的导出范畴》中的引理
\ref{spaces-perfect-lemma-flat-proper-perfect-direct-image-general}）。
上面证明的消失说明 $E \otimes^\mathbf{L} \kappa(s)$ 等于置于次数 $0$
的 $H^0(X_s, \omega_{X_s}^{\otimes m})$。缩小 $S$ 后可得，
$E = f_*\omega_{X/S}^{\otimes m}$ 是置于次数 $0$ 的局部自由
$\mathcal{O}_S$-模（而且如前所述，其构造与任意基变换交换）；见
《空间的导出范畴》中的引理
\ref{spaces-perfect-lemma-open-where-cohomology-in-degree-i-rank-r-geometric}。
态射 $f^*f_*\omega_{X/S}^{\otimes m} \to \omega_{X/S}^{\otimes m}$
限制到 $X_s$ 后满射。因此，它在 $X_s$ 的某个开邻域上满射。
由于 $f$ 是真态射，该开邻域包含 $f^{-1}(U)$，其中 $U$
是 $S$ 中 $s$ 的某个开邻域。

\medskip\noindent
假设 (1) 且 $g = 1$。由《代数曲线》中的引理
\ref{curves-lemma-contracting-rational-tails}，假设 (1) 意味着
$\omega_{X_s}$ 同构于 $\mathcal{O}_{X_s}$。若能证明缩小 $S$ 后
% P12-ERR-0281: see ../control/ERRATA_CANDIDATES.jsonl
可逆层 $\omega_{X/S}$ 变为平凡层，结论便成立。可以假设 $S$ 仿射。
进一步缩小 $S$ 后，可以写成
$$
Rf_*\mathcal{O}_X = (\mathcal{O}_S \xrightarrow{0} \mathcal{O}_S)
$$
它位于次数 $0$ 和 $1$，且与进一步的基变换相容；见引理
\ref{moduli-curves-lemma-genus}。由对偶性，这意味着
$$
Rf_*\omega_{X/S} = (\mathcal{O}_S \xrightarrow{0} \mathcal{O}_S)
$$
位于次数 $0$ 和 $1$\footnote{使用
$Rf_*\omega_{X/S}^\bullet =
% P12-ERR-0282: see ../control/ERRATA_CANDIDATES.jsonl
Rf_*R\SheafHom_{\mathcal{O}_X}(\mathcal{O}_X. \omega_{X/S}^\bullet) =
R\SheafHom_{\mathcal{O}_S}(Rf_*\mathcal{O}_X, \mathcal{O}_S)$；
这由《空间的对偶性》中的引理
\ref{spaces-duality-lemma-iso-on-RSheafHom} 和注
\ref{spaces-duality-remark-iso-on-RSheafHom} 得出；然后再使用由第
\ref{moduli-curves-section-relative-dualizing} 节中的定义得到的
$\omega_{X/S}^\bullet = \omega_{X/S}[1]$。}。
特别地，可以得到与基变换相容的同构
$\mathcal{O}_S \to f_*\omega_{X/S}$，因为
$Rf_*\omega_{X/S}$ 的构造与基变换相容（见上面给出的参考）。
由伴随性，得到整体截面 $\sigma \in \Gamma(X, \omega_{X/S})$。
该截面在纤维 $X_s$ 上的限制非零（事实上是一个基元素）；而且由于
$\omega_{X_s}$ 在该纤维上平凡，该截面在 $X_s$ 上处处非零。
% P12-ERR-0283: see ../control/ERRATA_CANDIDATES.jsonl
% P12-ERR-0284: see ../control/ERRATA_CANDIDATES.jsonl
因此，它在 $X_s$ 的某个开邻域上处处非零。由于 $f$ 是真态射，
该开邻域包含 $f^{-1}(U)$，其中 $U$ 是 $S$ 中 $s$ 的某个开邻域。
\end{proof}

\noindent
受引理 \ref{moduli-curves-lemma-semistable} 启发，我们作如下定义。

\begin{definition}
\label{moduli-curves-definition-semistable}
设 $f : X \to S$ 是曲线族。若
\begin{enumerate}
\item $X \to S$ 是预稳定曲线族；并且
\item 对所有 $s \in S$，$X_s$ 的亏格均为 $\geq 1$，且没有有理尾，
\end{enumerate}
则称 $f$ 是一个{\it 半稳定曲线族}。
\end{definition}

\noindent
特别地，亏格为 $0$ 的预稳定曲线族绝不半稳定。
设 $X$ 是域 $k$ 上的真概形，且 $\dim(X) \leq 1$。
则 $X \to \Spec(k)$ 是曲线族，因而可以问它是否半稳定。
展开定义可见，下列条件等价：
\begin{enumerate}
\item $X$ 半稳定；
\item $X$ 预稳定，亏格为 $\geq 1$，且没有有理尾；
\item $X_{\overline{k}}$ 连通，除有限多个结点外在 $\overline{k}$ 上光滑，
亏格为 $\geq 1$，并且没有同构于 $\mathbf{P}^1_{\overline{k}}$
且仅在一点与 $X_{\overline{k}}$ 的其余部分相交的不可约分支。
\end{enumerate}
要看出 (2) 与 (3) 等价，使用《代数曲线》中的引理
\ref{curves-lemma-contracting-rational-tails}：$X$ 没有有理尾，
当且仅当 $X_{\overline{k}}$ 没有有理尾。
这说明我们的定义与文献中的大多数定义一致。

\begin{lemma}
\label{moduli-curves-lemma-semistable-curves}
% P12-ERR-0285: see ../control/ERRATA_CANDIDATES.jsonl
存在开子叠 $\Curvesstack^{semistable} \subset \Curvesstack$，使得：
\begin{enumerate}
\item 给定曲线族 $f : X \to S$，下列条件等价：
\begin{enumerate}
\item 分类态射 $S \to \Curvesstack$ 经
$\Curvesstack^{semistable}$ 分解；
\item $X \to S$ 是半稳定曲线族；
\end{enumerate}
% P12-ERR-0286: see ../control/ERRATA_CANDIDATES.jsonl
\item 给定域 $k$ 上的真概形 $X$，且 $\dim(X) \leq 1$，
下列条件等价：
\begin{enumerate}
\item 分类态射 $\Spec(k) \to \Curvesstack$ 经
$\Curvesstack^{semistable}$ 分解；
\item $X$ 的奇点至多为结点，$\dim(X) = 1$，
$k = H^0(X, \mathcal{O}_X)$，$X$ 的亏格为 $\geq 1$，且
$X$ 没有有理尾；
\item $X$ 的奇点至多为结点，$\dim(X) = 1$，
$k = H^0(X, \mathcal{O}_X)$，而且对 $m \geq 2$，
% P12-ERR-0287: see ../control/ERRATA_CANDIDATES.jsonl
$\omega_{X_s}^{\otimes m}$ 整体生成。
\end{enumerate}
\end{enumerate}
\end{lemma}

\begin{proof}
(2)(b) 与 (2)(c) 的等价性是《代数曲线》中的引理
\ref{curves-lemma-contracting-rational-tails}。
在证明的其余部分，我们将按照定义 \ref{moduli-curves-definition-semistable} 使用 (2)(b)。

\medskip\noindent
由第 \ref{moduli-curves-section-open} 节的讨论，只需考虑预稳定曲线族
$f : X \to S$。由引理 \ref{moduli-curves-lemma-semistable}，得到所求轨迹的开性。
该开集的构造与任意基变换交换，因为由《代数曲线》中的引理
\ref{curves-lemma-contracting-rational-tails}，有理尾的存在或不存在
不受基域扩张影响。
\end{proof}

\begin{lemma}
\label{moduli-curves-lemma-semistable-one-piece-per-genus}
存在开闭子叠的分解
$$
\Curvesstack^{semistable} = \coprod\nolimits_{g \geq 1}
\Curvesstack^{semistable}_g
$$
其中每个 $\Curvesstack^{semistable}_g$ 由下列性质刻画：
\begin{enumerate}
\item 给定曲线族 $f : X \to S$，下列条件等价：
\begin{enumerate}
\item 分类态射 $S \to \Curvesstack$ 经
$\Curvesstack^{semistable}_g$ 分解；
\item $X \to S$ 是半稳定曲线族，且 $R^1f_*\mathcal{O}_X$
是秩为 $g$ 的局部自由 $\mathcal{O}_S$-模；
\end{enumerate}
% P12-ERR-0288: see ../control/ERRATA_CANDIDATES.jsonl
\item 给定域 $k$ 上的真概形 $X$，且 $\dim(X) \leq 1$，
下列条件等价：
\begin{enumerate}
\item 分类态射 $\Spec(k) \to \Curvesstack$ 经
$\Curvesstack^{semistable}_g$ 分解；
\item $X$ 的奇点至多为结点，$\dim(X) = 1$，
$k = H^0(X, \mathcal{O}_X)$，$X$ 的亏格为 $g$，且 $X$ 没有有理尾；
\item $X$ 的奇点至多为结点，$\dim(X) = 1$，
$k = H^0(X, \mathcal{O}_X)$，$X$ 的亏格为 $g$，而且对 $m \geq 2$，
% P12-ERR-0289: see ../control/ERRATA_CANDIDATES.jsonl
$\omega_{X_s}^{\otimes m}$ 整体生成。
\end{enumerate}
\end{enumerate}
\end{lemma}

\begin{proof}
合并引理 \ref{moduli-curves-lemma-semistable-curves} 和
\ref{moduli-curves-lemma-prestable-one-piece-per-genus} 即得。
\end{proof}

\begin{lemma}
\label{moduli-curves-lemma-semistable-curves-smooth}
态射 $\Curvesstack^{semistable} \to \Spec(\mathbf{Z})$ 和
$\Curvesstack^{semistable}_g \to \Spec(\mathbf{Z})$ 光滑。
\end{lemma}

\begin{proof}
由于 $\Curvesstack^{semistable}$ 是 $\Curvesstack^{nodal}$ 的开子叠，
结论由引理 \ref{moduli-curves-lemma-nodal-curves-smooth} 得出。
\end{proof}

\section{稳定曲线}
\label{moduli-curves-section-stable-curves}

\noindent
下述引理将帮助我们理解稳定曲线族。

\begin{lemma}
\label{moduli-curves-lemma-stable}
设 $f : X \to S$ 是亏格为 $g \geq 2$ 的预稳定曲线族，
并设 $s \in S$ 是基概形上的一点。下列条件等价：
\begin{enumerate}
\item $X_s$ 没有有理尾，也没有有理桥
（《代数曲线》中的例 \ref{curves-example-rational-tail} 和
\ref{curves-example-rational-bridge}）；
\item 对某个开子集 $s \in U \subset S$，$\omega_{X/S}$
在 $f^{-1}(U)$ 上丰沛。
\end{enumerate}
\end{lemma}

\begin{proof}
假设 (2)。则 $\omega_{X_s}$ 在 $X_s$ 上丰沛。
由《代数曲线》中的引理 \ref{curves-lemma-rational-tail-negative} 和
\ref{curves-lemma-rational-bridge-zero}，可知 (1) 成立
（还使用《簇》中的引理
\ref{varieties-lemma-ampleness-in-terms-of-degrees-components}
对丰沛可逆层的刻画）。

\medskip\noindent
假设 (1)。则由《代数曲线》中的引理
% P12-ERR-0290: see ../control/ERRATA_CANDIDATES.jsonl
\ref{curves-lemma-contracting-rational-bridges}，$\omega_{X_s}$
在 $X_s$ 上丰沛。结论由《空间的下降》中的引理
\ref{spaces-descent-lemma-ample-in-neighbourhood} 得出。
\end{proof}

\noindent
受引理 \ref{moduli-curves-lemma-stable} 启发，我们作如下定义。

\begin{definition}
\label{moduli-curves-definition-stable}
设 $f : X \to S$ 是曲线族。若
\begin{enumerate}
\item $X \to S$ 是预稳定曲线族；并且
% P12-ERR-0291: see ../control/ERRATA_CANDIDATES.jsonl
\item 对所有 $s \in S$，$X_s$ 的亏格均为 $\geq 2$，
且没有有理尾或有理桥，
\end{enumerate}
则称 $f$ 是一个{\it 稳定曲线族}。
\end{definition}

\noindent
特别地，亏格为 $0$ 或 $1$ 的预稳定曲线族绝不稳定。
设 $X$ 是域 $k$ 上的真概形，且 $\dim(X) \leq 1$。
则 $X \to \Spec(k)$ 是曲线族，因而可以问它是否稳定。
展开定义可见，下列条件等价：
\begin{enumerate}
\item $X$ 稳定；
\item $X$ 预稳定，亏格为 $\geq 2$，没有有理尾，也没有有理桥；
\item $X$ 几何连通，除有限多个结点外在 $k$ 上光滑，且
$\omega_X$ 丰沛。
\end{enumerate}
要看出 (2) 与 (3) 等价，使用上面的引理 \ref{moduli-curves-lemma-stable}。
这说明我们的定义与文献中的大多数定义一致。

\begin{lemma}
\label{moduli-curves-lemma-stable-curves}
% P12-ERR-0292: see ../control/ERRATA_CANDIDATES.jsonl
存在开子叠 $\Curvesstack^{stable} \subset \Curvesstack$，使得：
\begin{enumerate}
\item 给定曲线族 $f : X \to S$，下列条件等价：
\begin{enumerate}
\item 分类态射 $S \to \Curvesstack$ 经 $\Curvesstack^{stable}$ 分解；
\item $X \to S$ 是稳定曲线族；
\end{enumerate}
% P12-ERR-0293: see ../control/ERRATA_CANDIDATES.jsonl
\item 给定域 $k$ 上的真概形 $X$，且 $\dim(X) \leq 1$，
下列条件等价：
\begin{enumerate}
\item 分类态射 $\Spec(k) \to \Curvesstack$ 经
$\Curvesstack^{stable}$ 分解；
\item $X$ 的奇点至多为结点，$\dim(X) = 1$，
$k = H^0(X, \mathcal{O}_X)$，$X$ 的亏格为 $\geq 2$，
且 $X$ 没有有理尾或有理桥；
\item $X$ 的奇点至多为结点，$\dim(X) = 1$，
$k = H^0(X, \mathcal{O}_X)$，且
% P12-ERR-0294: see ../control/ERRATA_CANDIDATES.jsonl
$\omega_{X_s}$ 丰沛。
\end{enumerate}
\end{enumerate}
\end{lemma}

\begin{proof}
由第 \ref{moduli-curves-section-open} 节的讨论，只需考虑预稳定曲线族
$f : X \to S$。由引理 \ref{moduli-curves-lemma-stable}，得到所求轨迹的开性。
该开集的构造与任意基变换交换：一方面，由《代数曲线》中的引理
\ref{curves-lemma-contracting-rational-tails} 和
\ref{curves-lemma-contracting-rational-bridges}，有理尾或有理桥的
存在或不存在不受基域扩张影响；另一方面，由《下降》中的引理
\ref{descent-lemma-descending-property-ample}，丰沛性也不受基域扩张影响。
\end{proof}

\begin{definition}
\label{moduli-curves-definition-deligne-mumford}
\begin{reference}
\cite{DM}
\end{reference}
% P12-ERR-0295: see ../control/ERRATA_CANDIDATES.jsonl
我们用 $\overline{\mathcal{M}}$ 表示引理 \ref{moduli-curves-lemma-stable-curves}
中引入的、参数化稳定曲线族的代数叠 $\Curvesstack^{stable}$，
并称它为{\it 稳定曲线模叠}。对 $g \geq 2$，我们用
$\overline{\mathcal{M}}_g$ 表示引理
\ref{moduli-curves-lemma-stable-one-piece-per-genus} 中引入的代数叠，
并称它为{\it 亏格为 $g$ 的稳定曲线模叠}。
\end{definition}

\noindent
下面是必不可少的引理。

\begin{lemma}
\label{moduli-curves-lemma-stable-one-piece-per-genus}
存在开闭子叠的分解
$$
\overline{\mathcal{M}} = \coprod\nolimits_{g \geq 2} \overline{\mathcal{M}}_g
$$
其中每个 $\overline{\mathcal{M}}_g$ 由下列性质刻画：
\begin{enumerate}
\item 给定曲线族 $f : X \to S$，下列条件等价：
\begin{enumerate}
\item 分类态射 $S \to \Curvesstack$ 经
$\overline{\mathcal{M}}_g$ 分解；
\item $X \to S$ 是稳定曲线族，且 $R^1f_*\mathcal{O}_X$
是秩为 $g$ 的局部自由 $\mathcal{O}_S$-模；
\end{enumerate}
% P12-ERR-0296: see ../control/ERRATA_CANDIDATES.jsonl
\item 给定域 $k$ 上的真概形 $X$，且 $\dim(X) \leq 1$，
下列条件等价：
\begin{enumerate}
\item 分类态射 $\Spec(k) \to \Curvesstack$ 经
$\overline{\mathcal{M}}_g$ 分解；
\item $X$ 的奇点至多为结点，$\dim(X) = 1$，
$k = H^0(X, \mathcal{O}_X)$，$X$ 的亏格为 $g$，且
$X$ 没有有理尾或有理桥；
\item $X$ 的奇点至多为结点，$\dim(X) = 1$，
$k = H^0(X, \mathcal{O}_X)$，$X$ 的亏格为 $g$，且
% P12-ERR-0297: see ../control/ERRATA_CANDIDATES.jsonl
$\omega_{X_s}$ 丰沛。
\end{enumerate}
\end{enumerate}
\end{lemma}

\begin{proof}
合并引理 \ref{moduli-curves-lemma-stable-curves} 和
\ref{moduli-curves-lemma-prestable-one-piece-per-genus} 即得。
\end{proof}

\begin{lemma}
\label{moduli-curves-lemma-stable-curves-smooth}
态射 $\overline{\mathcal{M}} \to \Spec(\mathbf{Z})$ 和
$\overline{\mathcal{M}}_g \to \Spec(\mathbf{Z})$ 光滑。
\end{lemma}

\begin{proof}
由于 $\overline{\mathcal{M}}$ 是 $\Curvesstack^{nodal}$ 的开子叠，
结论由引理 \ref{moduli-curves-lemma-nodal-curves-smooth} 得出。
\end{proof}

\begin{lemma}
\label{moduli-curves-lemma-stable-curves-deligne-mumford}
叠 $\overline{\mathcal{M}}$ 和 $\overline{\mathcal{M}}_g$
是 $\Curvesstack^{DM}$ 的开子叠。特别地，
$\overline{\mathcal{M}}$ 和 $\overline{\mathcal{M}}_g$ 既是 DM 叠
（《叠的态射》中的定义
\ref{stacks-morphisms-definition-absolute-separated}），
也是 Deligne--Mumford 叠
（《代数叠》中的定义 \ref{algebraic-definition-deligne-mumford}）。
\end{lemma}

\begin{proof}
先证第一个断言。设 $X$ 是域 $k$ 上的真概形，其奇点至多为结点，
$\dim(X) = 1$，$k = H^0(X, \mathcal{O}_X)$，$X$ 的亏格为
$\geq 2$，而且 $X$ 没有有理尾或有理桥。需要证明分类态射
$\Spec(k) \to \overline{\mathcal{M}} \to \Curvesstack$
经 $\Curvesstack^{DM}$ 分解。可以先用代数闭包替换 $k$
（因为已经知道相关的叠是代数叠 $\Curvesstack$ 的开子叠）。
由引理 \ref{moduli-curves-lemma-stable-curves}、\ref{moduli-curves-lemma-DM-curves} 和
\ref{moduli-curves-lemma-in-DM-locus-vector-fields}，只需证明
$\text{Der}_k(\mathcal{O}_X, \mathcal{O}_X) = 0$。
这已在《代数曲线》中的引理
\ref{curves-lemma-stable-vector-fields} 中证明。

\medskip\noindent
由于 $\Curvesstack^{DM}$ 是 $\Curvesstack$ 中为 DM 的最大开子叠，
可知 $\Curvesstack^{DM}$ 的开子叠 $\overline{\mathcal{M}}$
也具有此性质。最后，由《叠的态射》中的定理
\ref{stacks-morphisms-theorem-DM}，DM 代数叠是 Deligne--Mumford 叠。
\end{proof}

\begin{lemma}
\label{moduli-curves-lemma-smooth-dense-in-stable}
设 $g \geq 2$。包含关系
$$
|\mathcal{M}_g| \subset |\overline{\mathcal{M}}_g|
$$
把左边实现为右边的开稠密子集。
\end{lemma}

\begin{proof}
由于 $\overline{\mathcal{M}}_g \subset \Curvesstack^{lci+}$ 是开的，
且 $\Curvesstack^{smooth} \cap \overline{\mathcal{M}}_g = \mathcal{M}_g$，
结论立即由引理 \ref{moduli-curves-lemma-smooth-dense} 得出。
\end{proof}

\section{收缩态射}
\label{moduli-curves-section-contracting}

\noindent
我们建议读者在继续本节之前，先熟悉《代数曲线》中的第
\ref{curves-section-contracting-rational-tails}、
\ref{curves-section-contracting-rational-bridges} 和
\ref{curves-section-contracting-to-stable} 节。本节的主要结果是“稳定化”态射
$$
\Curvesstack^{prestable}_g
\longrightarrow
\overline{\mathcal{M}}_g
$$
的存在性；见引理 \ref{moduli-curves-lemma-stabilization-morphism}。粗略地说，
该态射把亏格为 $g$ 的结点曲线的模点送到相伴稳定曲线的模点；
后者构造于《代数曲线》中的引理
\ref{curves-lemma-characterize-contraction-to-stable}。

\begin{lemma}
\label{moduli-curves-lemma-contract}
设 $S$ 是概形，$s \in S$ 是一点。设 $f : X \to S$ 和
$g : Y \to S$ 是曲线族，并设 $c : X \to Y$ 是 $S$ 上的态射。
若 $c_{s, *}\mathcal{O}_{X_s} = \mathcal{O}_{Y_s}$ 且
$R^1c_{s, *}\mathcal{O}_{X_s} = 0$，则以 $s$ 的一个开邻域替换 $S$ 后，
有 $\mathcal{O}_Y = c_*\mathcal{O}_X$ 和
$R^1c_*\mathcal{O}_X = 0$，而且经任意态射 $S' \to S$ 基变换后
这些等式仍成立。
\end{lemma}

\begin{proof}
设 $(U, u) \to (S, s)$ 是一个 \'etale 邻域，使得
$\mathcal{O}_{Y_U} = (X_U \to Y_U)_*\mathcal{O}_{X_U}$ 且
$R^1(X_U \to Y_U)_*\mathcal{O}_{X_U} = 0$，并且经任意
$U' \to U$ 基变换后也有相同结论。然后用 $U \to S$ 的开像替换 $S$。
给定 $S' \to S$，置 $U' = U \times_S S'$，便得到 \'etale 覆盖
$\{U' \to S'\}$ 和 $\{Y_{U'} \to Y_{S'}\}$。
因此，$c$ 经 $S' \to S$ 基变换后命题成立，可由
$X_U \to Y_U$ 经 $U' \to U$ 基变换后命题成立推出。
换言之，该问题在 $S$ 的 \'etale 拓扑下是局部的。
因此，由引理 \ref{moduli-curves-lemma-etale-locally-scheme}，可以假设 $X$ 和 $Y$
是概形。由《态射进阶》中的引理
\ref{more-morphisms-lemma-h1-fibre-zero-isom}，存在包含 $Y_s$
的开子概形 $V \subset Y$，使得
$c_*\mathcal{O}_X|_V = \mathcal{O}_V$ 且
$R^1c_*\mathcal{O}_X|_V = 0$，并且经任意 $S' \to S$ 基变换后
这些等式仍成立。由于 $g : Y \to S$ 是真态射，可以找到 $s$
的开邻域 $U \subset S$，使得 $g^{-1}(U) \subset V$。这个 $U$ 即可。
\end{proof}

\begin{lemma}
\label{moduli-curves-lemma-contract-basic-uniqueness}
设 $S$ 是概形，$s \in S$ 是一点。设 $f : X \to S$ 和
$g_i : Y_i \to S$（$i = 1, 2$）是曲线族，并设
$c_i : X \to Y_i$ 是 $S$ 上的态射。假设存在与 $c_{1, s}$ 和
$c_{2, s}$ 相容的纤维同构 $Y_{1, s} \cong Y_{2, s}$。
若 $c_{1, s, *}\mathcal{O}_{X_s} = \mathcal{O}_{Y_{1, s}}$ 且
$R^1c_{1, s, *}\mathcal{O}_{X_s} = 0$，则存在 $s$ 的开邻域 $U$
以及 $U$ 上曲线族的同构 $Y_{1, U} \cong Y_{2, U}$；该同构与给定的
纤维同构以及 $c_1$、$c_2$ 相容。
\end{lemma}

\begin{proof}
回顾 $\mathcal{O}_{S, s} = \colim \mathcal{O}_S(U)$，其中余极限取遍
$s$ 的仿射邻域 $U$ 所成的系统。因此，局部环上有限表示代数空间的范畴
是 $s$ 的仿射邻域上有限表示代数空间范畴的余极限。见《空间的极限》中的引理
\ref{spaces-limits-lemma-descend-finite-presentation}。
这样便归约到 $S$ 是某个局部环的谱且 $s$ 是闭点的情形。

\medskip\noindent
假设 $S = \Spec(A)$，其中 $A$ 是局部环，$s$ 是闭点。
将 $A = \colim A_j$ 写成局部 Noether 环 $A_j$ 的余极限
（例如可取在 $\mathbf{Z}$ 上本质有限型），其过渡同态均为局部同态。
置 $S_j = \Spec(A_j)$，其闭点为 $s_j$。可以找到一个 $j$ 以及曲线族
$X_j \to S_j$、$Y_{j, i} \to S_j$；见引理
\ref{moduli-curves-lemma-curves-qs-lfp} 和《叠的极限》中的引理
\ref{stacks-limits-lemma-representable-by-spaces-limit-preserving}。
必要时增大 $j$，可以找到态射 $c_{j, i} : X_j \to Y_{j, i}$，
其到 $s$ 的基变换为 $c_i$；见《空间的极限》中的引理
\ref{spaces-limits-lemma-descend-finite-presentation}。
由于 $\kappa(s) = \colim \kappa(s_j)$，同样可以假设存在与
$c_{j, 1, s_j}$ 和 $c_{j, 2, s_j}$ 相容的同构
$Y_{j, 1, s_j} \cong Y_{j, 2, s_j}$。最后，假设
$c_{1, s, *}\mathcal{O}_{X_s} = \mathcal{O}_{Y_{1, s}}$ 和
$R^1c_{1, s, *}\mathcal{O}_{X_s} = 0$ 可由 $c_{j, 1, s_j}$ 继承，因为
% P12-ERR-0298: see ../control/ERRATA_CANDIDATES.jsonl
$\{s_j \to s\}$ 是 fpqc 覆盖，而 $c_{1, s}$ 是
$c_{j, 1, s_j}$ 经该覆盖的基变换（细节略去）。
这样便把引理归约到下一段所讨论的情形。

\medskip\noindent
假设 $S$ 是 Noether 局部环 $\Lambda$ 的谱，$s$ 是闭点。
考虑态射
$$
(c_1, c_2) : X \longrightarrow Y_1 \times_S Y_2
$$
的概形论像 $Z$。引理的陈述等价于断言：$Z$ 经投影态射同构地映到
$Y_1$ 和 $Y_2$。由于取该态射的概形论像与平坦基变换交换
% P12-ERR-0299: see ../control/ERRATA_CANDIDATES.jsonl
（《空间的态射》中的引理
\ref{spaces-morphisms-lemma-flat-base-change-scheme-theoretic-image}，
可以用 $\Lambda$ 的完备化替换它
（《代数进阶》中的第 \ref{more-algebra-section-permanence-completion} 节）。

\medskip\noindent
假设 $S$ 是完备 Noether 局部环 $\Lambda$ 的谱。
注意，此时 $X$、$Y_1$、$Y_2$ 都是概形
（《空间的态射进阶》中的引理
\ref{spaces-more-morphisms-lemma-projective-over-complete}）。
% P12-ERR-0300: see ../control/ERRATA_CANDIDATES.jsonl
用 $X_n$、$Y_{1, n}$、$Y_{2, n}$ 表示 $X$、$Y_1$、$Y_2$
到 $\Spec(\Lambda/\mathfrak m^{n + 1})$ 的基变换。回顾箭头
$$
\Deformationcategory_{X_s \to Y_{2, s}} \cong
\Deformationcategory_{X_s \to Y_{1, s}} \longrightarrow
\Deformationcategory_{X_s}
$$
是等价；见《形变问题》中的引理
\ref{examples-defos-lemma-schemes-morphisms-smooth-to-base}。
因此，在 $\Deformationcategory_{X_s \to Y_{1, s}}$ 中存在形式对象的同构
$(X_n \to Y_{1, n}) \cong (X_n \to Y_{2, n})$。
最后，由 Grothendieck 代数化定理（《概形的上同调》中的引理
\ref{coherent-lemma-algebraize-morphism}），这给出与 $c_1$ 和 $c_2$
相容的同构 $Y_1 \to Y_2$。
\end{proof}

\begin{lemma}
\label{moduli-curves-lemma-contract-basic}
设 $f : X \to S$ 是曲线族，$s \in S$ 是一点。设
$h_0 : X_s \to Y_0$ 是到 $\kappa(s)$ 上真概形 $Y_0$ 的态射，
满足 $h_{0, *}\mathcal{O}_{X_s} = \mathcal{O}_{Y_0}$ 且
$R^1h_{0, *}\mathcal{O}_{X_s} = 0$。则存在基本 \'etale 邻域
$(U, u) \to (S, s)$、曲线族 $Y \to U$ 以及 $U$ 上的态射
$h : X_U \to Y$，使得它在 $u$ 上的纤维同构于 $h_0$。
\end{lemma}

\begin{proof}
先作一些归约；建议读者跳到后面。该问题在 $S$ 上是局部的，
故可以假设 $S$ 仿射。将 $S = \lim S_i$ 写成在 $\mathbf{Z}$ 上
有限型仿射概形 $S_i$ 的余滤极限。对某个 $i$，可以找到曲线族
$X_i \to S_i$，其基变换为 $X \to S$。这由引理
\ref{moduli-curves-lemma-curves-qs-lfp} 和《叠的极限》中的引理
\ref{stacks-limits-lemma-representable-by-spaces-limit-preserving} 得出。
设 $s_i \in S_i$ 是 $s$ 的像。注意
$\kappa(s) = \colim \kappa(s_i)$，而且 $X_s$ 是概形
（《域上的空间》中的引理
\ref{spaces-over-fields-lemma-codim-1-point-in-schematic-locus}）。
增大 $i$ 后，可以假设存在 $\kappa(s_i)$ 上有限型概形之间的态射
$h_{i, 0} : X_{i, s_i} \to Y_i$，其到 $\kappa(s)$ 的基变换为 $h_0$；
见《极限》中的引理 \ref{limits-lemma-descend-finite-presentation}。
再次增大 $i$ 后，可以假设 $Y_i$ 在 $\kappa(s_i)$ 上为真；
见《极限》中的引理 \ref{limits-lemma-eventually-proper}。
% P12-ERR-0301: see ../control/ERRATA_CANDIDATES.jsonl
设 $g_{i, 0} : Y_0 \to Y_{i, 0}$ 为投影。注意，它作为
$\Spec(\kappa(s)) \to \Spec(\kappa(s_i))$ 的基变换，是忠实平坦态射。
由平坦基变换，有
$$
h_{0, *}\mathcal{O}_{X_s} = g_{i, 0}^*h_{i, 0, *}\mathcal{O}_{X_{i, s_i}}
\quad\text{且}\quad
% P12-ERR-0302: see ../control/ERRATA_CANDIDATES.jsonl
R^1h_{0, *}\mathcal{O}_{X_s} = g_{i, 0}^*Rh_{i, 0, *}\mathcal{O}_{X_{i, s_i}}
$$
见《概形的上同调》中的引理
\ref{coherent-lemma-flat-base-change-cohomology}。
由忠实平坦性可知，$X_i \to S_i$、$s_i \in S_i$ 和
% P12-ERR-0303: see ../control/ERRATA_CANDIDATES.jsonl
$X_{i, s_i} \to Y_i$ 满足该引理的全部假设。
这便归约到下一段所讨论的情形。

\medskip\noindent
假设 $S$ 在 $\mathbf{Z}$ 上仿射且有限型。设
$\mathcal{O}_{S, s}^h$ 是 $S$ 在 $s$ 处局部环的 Hensel 化。
由《代数进阶》中的引理 \ref{more-algebra-lemma-henselization-G-ring}
和命题 \ref{more-algebra-proposition-ubiquity-G-ring}，
$\mathcal{O}_{S, s}^h$ 是 G-环。假设能够构造曲线族
$Y' \to \Spec(\mathcal{O}_{S, s}^h)$ 以及
$\Spec(\mathcal{O}_{S, s}^h)$ 上的态射
$$
h' : X \times_S \Spec(\mathcal{O}_{S, s}^h) \longrightarrow Y'
$$
其到闭点的基变换为 $h_0$。这已经足够。事实上，首先使用
$$
\mathcal{O}_{S, s}^h = \colim_{(U, u)} \mathcal{O}_U(U)
$$
其中余极限取遍基本 \'etale 邻域所成的滤过范畴
（《态射进阶》中的引理
\ref{more-morphisms-lemma-describe-henselization}）。
其次，再次使用如下事实：给定 $Y'$ 后，可以把它下降为某个 $U$ 上的
$Y \to U$（见上面给出的参考）。然后使用《极限》中的引理
\ref{limits-lemma-descend-finite-presentation}，把 $h'$ 下降为某个 $h$。
这便归约到下一段所讨论的情形。

\medskip\noindent
假设 $S = \Spec(\Lambda)$，其中
$(\Lambda, \mathfrak m, \kappa)$ 是 Hensel Noether 局部 G-环，
而 $s$ 是 $S$ 的闭点。回顾映射
$$
\Deformationcategory_{X_s \to Y_0} \to \Deformationcategory_{X_s}
$$
是等价；见《形变问题》中的引理
\ref{examples-defos-lemma-schemes-morphisms-smooth-to-base}。
（这是证明中唯一重要的一步；其余全是技巧。）
% P12-ERR-0304: see ../control/ERRATA_CANDIDATES.jsonl
用 $\Lambda^\wedge$ 表示 $\mathfrak m$-进完备化。
$X$ 到 $\Lambda/\mathfrak m^{n + 1}$ 的拉回 $X_n$
在 $\Lambda^\wedge$ 上定义了 $\Deformationcategory_{X_s}$ 的形式对象
$\xi$。由上述等价，得到 $\Lambda^\wedge$ 上
$\Deformationcategory_{X_s \to Y_0}$ 的形式对象 $\xi'$。
因此得到一个巨大的交换图
$$
\xymatrix{
\ldots \ar[r] &
X_n \ar[r] \ar[d] &
X_{n - 1} \ar[r] \ar[d] &
\ldots \ar[r] &
X_s \ar[d] \\
\ldots \ar[r] &
Y_n \ar[r] \ar[d] &
Y_{n - 1} \ar[r] \ar[d] &
\ldots \ar[r] &
Y_0 \ar[d] \\
\ldots \ar[r] &
\Spec(\Lambda/\mathfrak m^{n + 1}) \ar[r] &
\Spec(\Lambda/\mathfrak m^n) \ar[r] &
\ldots \ar[r] &
\Spec(\kappa)
}
$$
由《Quot》中的引理 \ref{quot-lemma-curves-existence}，
形式对象 $(Y_n)$ 来自曲线族 $Y' \to \Spec(\Lambda^\wedge)$。
由《空间的态射进阶》中的引理
\ref{spaces-more-morphisms-lemma-algebraize-morphism}，得到态射
$h' : X_{\Lambda^\wedge} \to Y'$，它对所有 $n$ 都诱导给定态射
$X_n \to Y_n$，特别地诱导给定态射 $X_s \to Y_0$。

\medskip\noindent
最后作一个标准的代数化/逼近论证。首先注意，可以找到有限生成
$\Lambda$-子代数 $\Lambda \subset A \subset \Lambda^\wedge$、
曲线族 $Y'' \to \Spec(A)$ 以及 $A$ 上的态射
$h'' : X_A \to Y''$，其到 $\Lambda^\wedge$ 的基变换为 $h'$。
这是因为 $\Lambda^\wedge$ 是这些环 $A$ 的滤过余极限；可以像前面那样，
利用 $\Curvesstack$ 局部有限表示这一事实（由《叠的极限》中的引理
\ref{stacks-limits-lemma-representable-by-spaces-limit-preserving}
得到 $A$ 上的 $Y''$），并使用《空间的极限》中的引理
\ref{spaces-limits-lemma-descend-finite-presentation}，
把 $h'$ 下降为某个 $h''$。然后，可以使用 G-环的逼近性质
（采用《环同态的光滑化》中的定理
\ref{smoothing-theorem-approximation-property} 的形式），
找到映射 $A \to \Lambda$，它所诱导的 $A \to \kappa$
与 $A \to \Lambda^\wedge$ 所诱导的映射相同。
把 $h''$ 基变换到 $\Lambda$，证明即告完成。
\end{proof}

\begin{lemma}
\label{moduli-curves-lemma-contract-prestable-to-stable}
设 $f : X \to S$ 是亏格为 $g \geq 2$ 的预稳定曲线族。
$f$ 存在分解 $X \to Y \to S$，其中
% P12-ERR-0305: see ../control/ERRATA_CANDIDATES.jsonl
$g : Y \to S$ 是稳定曲线族，而 $c : X \to Y$ 具有下列性质：
\begin{enumerate}
\item $\mathcal{O}_Y = c_*\mathcal{O}_X$ 且
$R^1c_*\mathcal{O}_X = 0$，并且经任意态射 $S' \to S$
基变换后这些等式仍成立；
\item 对任意 $s \in S$，态射 $c_s : X_s \to Y_s$ 是
《代数曲线》中的第 \ref{curves-section-contracting-to-stable} 节
所讨论的有理尾和有理桥的收缩。
\end{enumerate}
此外，$c : X \to Y$ 在唯一同构意义下唯一。
\end{lemma}

\begin{proof}
设 $s \in S$。令 $c_0 : X_s \to Y_0$ 为《代数曲线》中的第
\ref{curves-section-contracting-to-stable} 节的收缩
（更精确地说，见《代数曲线》中的引理
\ref{curves-lemma-characterize-contraction-to-stable}）。
由引理 \ref{moduli-curves-lemma-contract-basic}，存在基本 \'etale 邻域
$(U, u)$ 以及 $U$ 上曲线族之间的态射 $c : X_U \to Y$，
其在 $u$ 处的纤维恢复 $c_0$。由于 $\omega_{Y_0}$ 丰沛，
必要时缩小 $U$ 后，由引理 \ref{moduli-curves-lemma-stable-curves} 和
\ref{moduli-curves-lemma-stable-one-piece-per-genus} 所蕴含的开性可知，
$Y \to U$ 是亏格为 $g$ 的稳定曲线族。必要时再次缩小 $U$ 后，
引理 \ref{moduli-curves-lemma-contract} 给出 $c : X_U \to Y$ 的断言 (1)。
此外，(2) 由《代数曲线》中的引理
\ref{curves-lemma-characterize-contraction-to-stable} 的唯一性成立。
由此可知，引理中的态射 $c$ 在 $S$ 上 \'etale 局部存在。
更精确地，存在 \'etale 覆盖 $\{U_i \to S\}$ 以及 $U_i$ 上的态射
$c_i : X_{U_i} \to Y_i$，其中 $Y_i \to U_i$ 是具有引理中性质
(1) 和 (2) 的稳定曲线族。

\medskip\noindent
为完成证明，只需证明 $c : X \to Y$ 的唯一性（在唯一同构意义下）。
事实上，一旦完成这一步，便得到同构
% P12-ERR-0306: see ../control/ERRATA_CANDIDATES.jsonl
$$
\varphi_{ij} :
Y_i \times_{U_i} (U_i \times_S U_j)
\longrightarrow
Y_i \times_{U_j} (U_i \times_S U_j)
$$
% P12-ERR-0307: see ../control/ERRATA_CANDIDATES.jsonl
它们在 $U_i \times U_j \times U_k$ 上满足上链条件（由唯一性）。
% P12-ERR-0308: see ../control/ERRATA_CANDIDATES.jsonl
由于 $\overline{\mathcal{M}_g}$ 是代数叠，下降数据有效，
从而得到 $Y \to S$。态射 $c_i$ 下降为 $S$ 上的态射
$c : X \to Y$。最后，$c$ 的性质 (1) 和 (2) 立即由 $c_i$
的性质 (1) 和 (2) 得出。

\medskip\noindent
最后，若
% P12-ERR-0309: see ../control/ERRATA_CANDIDATES.jsonl
$c_1 : X \to Y_i$（$i = 1, 2$）是到 $S$ 上两个
% P12-ERR-0310: see ../control/ERRATA_CANDIDATES.jsonl
稳定曲线族的态射，并满足 (1) 和 (2)，则由引理
\ref{moduli-curves-lemma-contract-basic-uniqueness}，至少在 $S$ 上局部，
得到与 $c_1$ 和 $c_2$ 相容的态射 $Y_1 \to Y_2$。
略去验证这些态射唯一的过程（提示：这由 $c_1$ 的概形论像是
$Y_1$ 这一事实得出）。因此，这些局部给出的态射可以粘合，证明完成。
\end{proof}

\begin{lemma}
\label{moduli-curves-lemma-stabilization-morphism}
设 $g \geq 2$。存在 $\mathbf{Z}$ 上代数叠之间的态射
% P12-ERR-0311: see ../control/ERRATA_CANDIDATES.jsonl
$$
stabilization :
\Curvesstack^{prestable}_g
\longrightarrow
\overline{\mathcal{M}}_g
$$
它把亏格为 $g$ 的预稳定曲线族 $X \to S$ 送到引理
\ref{moduli-curves-lemma-contract-prestable-to-stable} 中与之相伴的稳定曲线族
$Y \to S$。
\end{lemma}

\begin{proof}
为说明这一点，只需检验引理
\ref{moduli-curves-lemma-contract-prestable-to-stable} 的构造与基变换相容
（也与同构相容，但这是立即的）；见《叠的性质》中的第
\ref{stacks-properties-section-conventions} 节所引入的关于代数叠的
（语言滥用）约定。为此，只需检验引理
\ref{moduli-curves-lemma-contract-prestable-to-stable} 的性质 (1) 和 (2)
在基变换下稳定。对 (1)，这立即清楚。对 (2)，一方面，
《代数曲线》中的引理 \ref{curves-lemma-contracting-rational-tails} 和
\ref{curves-lemma-contracting-rational-bridges} 的收缩在基域扩张下稳定；
另一方面，《代数曲线》中的引理
\ref{curves-lemma-characterize-contraction-to-stable}
对纤维上态射的刻画条件在基域扩张下保持。因此 (2) 成立。
\end{proof}

\section{稳定约化定理}
\label{moduli-curves-section-stable-reduction}

\noindent
在半稳定约化一章中，我们已经证明了著名的曲线半稳定约化定理。
设 $K$ 是离散赋值环 $R$ 的分式域。设 $C$ 是 $K$ 上的射影光滑曲线，
满足 $K = H^0(C, \mathcal{O}_C)$。按照《半稳定约化》中的定义
\ref{models-definition-semistable}，若 $R$ 上存在泛纤维为 $C$
的预稳定曲线族，或者 $C$ 在 $R$ 上的某个（等价地，任意）
极小正则模型预稳定，则称 $C$ 具有{\it 半稳定约化}。
本节将说明，对亏格为 $g \geq 2$ 的曲线，这也等价于稳定约化。

\begin{lemma}
\label{moduli-curves-lemma-stable-reduction}
设 $R$ 是分式域为 $K$ 的离散赋值环。设 $C$ 是 $K$ 上亏格为
$g \geq 2$ 的光滑射影曲线，满足 $K = H^0(C, \mathcal{O}_C)$。
下列条件等价：
\begin{enumerate}
\item $C$ 具有半稳定约化
（《半稳定约化》中的定义 \ref{models-definition-semistable}）；
\item $R$ 上存在泛纤维为 $C$ 的稳定曲线族。
\end{enumerate}
\end{lemma}

\begin{proof}
稳定曲线族也是预稳定曲线族，故 (2) 蕴含 (1) 是立即的。
反过来，给定 $R$ 上泛纤维为 $C$ 的预稳定曲线族，
由引理 \ref{moduli-curves-lemma-contract-prestable-to-stable}，可以把它收缩成稳定曲线族。
由于泛纤维已经稳定，该过程不改变泛纤维，证明完成。
\end{proof}

\noindent
下述引理说明，引理 \ref{moduli-curves-lemma-stable-reduction} 所保证的 $R$
上稳定曲线族在唯一同构意义下唯一。

\begin{lemma}
\label{moduli-curves-lemma-unique-stable-model}
设 $R$ 是分式域为 $K$ 的离散赋值环。设 $C$ 是 $K$ 上亏格为 $g$
的光滑真曲线，满足 $K = H^0(C, \mathcal{O}_C)$。
若 $X$ 和 $X'$ 是 $C$ 的模型
（《半稳定约化》中的第 \ref{models-section-models} 节），
而且 $X$ 和 $X'$ 是 $R$ 上亏格为 $g$ 的稳定曲线族，
则存在唯一的模型同构 $X \to X'$。
\end{lemma}

\begin{proof}
设 $Y$ 是 $C$ 的极小模型。回顾 $Y$ 存在且唯一，
并且在 $R$ 上是相对维数为 $1$、至多有结点奇性的态射；
见《半稳定约化》中的命题
\ref{models-proposition-exists-minimal-model} 以及引理
\ref{models-lemma-minimal-model-unique} 和
\ref{models-lemma-semistable}（后者适用是因为我们有 $X$）。
存在收缩态射
$$
Y \longrightarrow Z
$$
使得 $Z$ 是 $R$ 上亏格为 $g$ 的稳定曲线族
（引理 \ref{moduli-curves-lemma-contract-prestable-to-stable}）。我们断言，
存在唯一的模型同构 $X \to Z$。由对称性，对 $X'$ 也有相同结论，
从而完成证明。

\medskip\noindent
由《半稳定约化》中的引理
\ref{models-lemma-blowup-at-worst-nodal}，存在序列
$$
X_m \to \ldots \to X_1 \to X_0 = X
$$
使得 $X_{i + 1} \to X_i$ 是 $X_i$ 的奇异闭点 $x_i$ 的吹起，
$X_i \to \Spec(R)$ 是相对维数为 $1$、至多有结点奇性的态射，
而 $X_m$ 正则。由《半稳定约化》中的引理
\ref{models-lemma-pre-exists-minimal-model}，存在 $C$ 的真正则模型序列
$$
X_m = Y_n \to Y_{n - 1} \to \ldots \to Y_1 \to Y_0 = Y
$$
使得每个态射都是第一类例外曲线的收缩\footnote{事实上，
有 $X_m = Y$；即 $X_m$ 不含任何第一类例外曲线。
我们鼓励读者仔细思考这一点，因为它会稍微简化证明。}。
由《半稳定约化》中的引理
\ref{models-lemma-blowdown-at-worst-nodal}，每个 $Y_i$
在 $R$ 上都是相对维数为 $1$、至多有结点奇性的态射。
为证明该断言，只需证明存在同构 $X \to Z$，它与态射
$X_m \to X$ 和 $X_m = Y_n \to Y \to Z$ 相容。
设 $s \in \Spec(R)$ 是闭点。由引理
\ref{moduli-curves-lemma-contract-basic-uniqueness} 或引理
\ref{moduli-curves-lemma-contract-prestable-to-stable}，问题归约为证明态射
$X_{m, s} \to X_s$ 和 $X_{m, s} \to Z_s$
都等于《代数曲线》中的引理
\ref{curves-lemma-characterize-contraction-to-stable} 的典范态射。

\medskip\noindent
对 $\kappa(s)$ 上概形之间的态射 $c : U \to V$，若对
$v \in V$ 有 $\dim(U_v) \leq 1$，并且
$\mathcal{O}_V = c_*\mathcal{O}_U$、$R^1c_*\mathcal{O}_U = 0$，
则称 $c$ 具有性质 (*)。该性质在复合下稳定。
由于 $X_s$ 和 $Z_s$ 都是 $\kappa(s)$ 上亏格为 $g$ 的稳定曲线，
只需证明态射 $Y_s \to Z_s$、$X_{i + 1, s} \to X_{i, s}$ 和
% P12-ERR-0312: see ../control/ERRATA_CANDIDATES.jsonl
$Y_{i + 1, s} \to Y_{i, s}$ 中的每一个都满足性质 (*)；见
《代数曲线》中的引理
\ref{curves-lemma-characterize-contraction-to-stable}。

\medskip\noindent
性质 (*) 对 $Y_s \to Z_s$ 由构造成立。

\medskip\noindent
态射 $c : X_{i + 1, s} \to X_{i, s}$ 构造并研究于
《半稳定约化》中的引理
\ref{models-lemma-blowup-at-worst-nodal} 的证明中。
只需在 $X_{i, s}$ 上 \'etale 局部检验 (*)。
因而只需对《半稳定约化》中的例 \ref{models-example-blowup}
的态射“$X_1 \to X_0$”到 $R/\pi R$ 的基变换检验 (*)。
把具体计算留给读者。

\medskip\noindent
态射 $c : Y_{i + 1, s} \to Y_{i, s}$ 是第一类例外曲线
$E \subset Y_{i + 1}$ 的吹下的限制；也就是说，
$b : Y_{i + 1} \to Y_i$ 是 $E$ 的收缩，也就是说，$b$ 是曲面
$Y_i$ 上一个正则点的吹起
（《曲面的解消》中的第 \ref{resolve-section-minus-one} 节）。
于是 $\mathcal{O}_{Y_i} = b_*\mathcal{O}_{Y_{i + 1}}$ 且
$R^1b_*\mathcal{O}_{Y_{i + 1}} = 0$；例如见《曲面的解消》中的引理
\ref{resolve-lemma-cohomology-of-blowup}。
由《态射进阶》中的引理
\ref{more-morphisms-lemma-check-h1-fibre-zero}、
\ref{more-morphisms-lemma-h1-fibre-zero} 和
\ref{more-morphisms-lemma-h1-fibre-zero-check-h0-kappa}，可知
$\mathcal{O}_{Y_{i, s}} = c_*\mathcal{O}_{Y_{i + 1, s}}$ 且
$R^1c_*\mathcal{O}_{Y_{i + 1, s}} = 0$
（最后一个引理只给出
$\mathcal{O}_{Y_{i, s}} \to c_*\mathcal{O}_{Y_{i + 1, s}}$
的满射性，但单射性容易由如下事实得出：$Y_{i, s}$ 约化，
而 $c$ 只改变一个闭点上方的部分）。证明完成。
\end{proof}

\noindent
由引理 \ref{moduli-curves-lemma-stable-reduction} 和《半稳定约化》中的定理
\ref{models-theorem-semistable-reduction}，立即得到稳定约化定理。

\begin{theorem}
\label{moduli-curves-theorem-stable-reduction}
\begin{reference}
\cite[Corollary 2.7]{DM}
\end{reference}
设 $R$ 是分式域为 $K$ 的离散赋值环。设 $C$ 是 $K$ 上亏格为
$g \geq 2$ 的光滑射影曲线，满足 $H^0(C, \mathcal{O}_C) = K$。则：
\begin{enumerate}
\item 存在离散赋值环扩张 $R \subset R'$，它诱导分式域的有限可分扩张
$K'/K$；并且存在该环上亏格为 $g$ 的稳定曲线族
$Y \to \Spec(R')$，满足在 $K'$ 上
$Y_{K'} \cong C_{K'}$；
\item 存在有限可分扩张 $L/K$ 以及亏格为 $g$ 的稳定曲线族
$Y \to \Spec(A)$，其中 $A \subset L$ 是 $R$ 在 $L$ 中的整闭包，
并且在 $L$ 上有 $Y_L \cong C_L$。
\end{enumerate}
\end{theorem}

\begin{proof}
(1) 立即由引理 \ref{moduli-curves-lemma-stable-reduction} 和《半稳定约化》中的定理
\ref{models-theorem-semistable-reduction} 得出。

\medskip\noindent
证明 (2)。设 $L/K$ 是《半稳定约化》中的定理
\ref{models-theorem-semistable-reduction} 第 (3) 部分所得的有限可分扩张。
设 $A \subset L$ 是 $R$ 的整闭包。回顾 $A$ 是 $R$ 上有限的
Dedekind 整环，具有有限多个极大理想
$\mathfrak m_1, \ldots, \mathfrak m_n$；见《代数进阶》中的注
\ref{more-algebra-remark-finite-separable-extension}。
置 $S = \Spec(A)$、$S_i = \Spec(A_{\mathfrak m_i})$、
$U = \Spec(L)$ 和 $U_i = S_i \setminus \{\mathfrak m_i\}$。
注意，对 $i = 1, \ldots, n$，有 $U \cong U_i$。
置 $X = C_L$，把它视为 $S$ 的开子概形 $U$ 上的概形。
由 $L$、$A$ 的选取以及引理 \ref{moduli-curves-lemma-stable-reduction}，
存在稳定曲线族 $X_i \to S_i$ 以及同构
$X \times_U U_i \cong X_i \times_{S_i} U_i$。
由《空间的极限》中的引理
\ref{spaces-limits-lemma-glueing-near-multiple-closed-points}，
可以找到有限表示态射 $Y \to S$，它到每个 $S_i$ 的基变换
都同构于 $X_i$（$i = 1, \ldots, n$）。或者，可以使用如下事实：
$S = \bigcup_{i = 1, \ldots, n} S_i$ 是 $S$ 的开覆盖，且当
$i \not = j$ 时 $S_i \cap S_j = U$；再应用 $n - 1$ 次
《空间的极限》中的引理 \ref{spaces-limits-lemma-relative-glueing}，
得到 $Y \to S$，它到每个 $S_i$ 的基变换都同构于 $X_i$
（$i = 1, \ldots, n$）。显然，$Y \to S$ 正是所求的稳定曲线族。
\end{proof}

\section{稳定曲线叠的性质}
\label{moduli-curves-section-properties-stable}

\noindent
本节证明 $g \geq 2$ 时 $\overline{\mathcal{M}}_g$ 的基本结构结果。

\begin{lemma}
\label{moduli-curves-lemma-stable-separated}
设 $g \geq 2$。叠 $\overline{\mathcal{M}}_g$ 分离。
\end{lemma}

\begin{proof}
该陈述意味着态射
$\overline{\mathcal{M}}_g \to \Spec(\mathbf{Z})$ 分离。
我们将使用《叠的态射进阶》中的引理
\ref{stacks-more-morphisms-lemma-refined-valuative-criterion-separated}
所陈述的精细 Noether 赋值判据来证明这一点。

\medskip\noindent
由于 $\overline{\mathcal{M}}_g$ 是 $\Curvesstack$ 的开子叠，
由引理 \ref{moduli-curves-lemma-curves-qs-lfp} 可知，
$\overline{\mathcal{M}}_g \to \Spec(\mathbf{Z})$ 拟分离且局部有限表示。
特别地，叠 $\overline{\mathcal{M}}_g$ 局部 Noether
（《叠的态射》中的引理
\ref{stacks-morphisms-lemma-locally-finite-type-locally-noetherian}）。
由引理 \ref{moduli-curves-lemma-smooth-dense-in-stable}，开浸入
$\mathcal{M}_g \to \overline{\mathcal{M}}_g$ 的像稠密。
此外，$\mathcal{M}_g \to \overline{\mathcal{M}}_g$ 拟紧
（《叠的态射》中的引理
\ref{stacks-morphisms-lemma-locally-closed-in-noetherian}），
因而是有限型的。因此，《叠的态射进阶》中的引理
\ref{stacks-more-morphisms-lemma-refined-valuative-criterion-separated}
对态射
$$
\mathcal{M}_g \to \overline{\mathcal{M}}_g
\quad\text{且}\quad
\overline{\mathcal{M}}_g \to \Spec(\mathbf{Z})
$$
的所有预备假设均成立，只需检验以下条件：给定任意 $2$-交换图
$$
\xymatrix{
\Spec(K) \ar[r] \ar[d] &
\mathcal{M}_g \ar[r] &
\overline{\mathcal{M}}_g \ar[d] \\
\Spec(R) \ar[rr] \ar@{..>}[rru] & & \Spec(\mathbf{Z})
}
$$
其中 $R$ 是分式域为 $K$ 的离散赋值环，则虚线箭头范畴或者为空，
或者是恰有一个同构类的 setoid。（注意，无需过多担心 $2$-箭头；
见《叠的态射》中的引理
\ref{stacks-morphisms-lemma-cat-dotted-arrows-independent}。）
展开其含义，并使用如下事实：$\mathcal{M}_g$ 和
$\overline{\mathcal{M}}_g$ 分别是参数化亏格为 $g$ 的光滑曲线族和
稳定曲线族的代数叠，便发现所需证明的正是引理
\ref{moduli-curves-lemma-unique-stable-model} 中陈述并证明的唯一性结果。
\end{proof}

\begin{lemma}
\label{moduli-curves-lemma-stable-quasi-compact}
设 $g \geq 2$。叠 $\overline{\mathcal{M}}_g$ 拟紧。
\end{lemma}

\begin{proof}
我们使用第 \ref{moduli-curves-section-polarized-curves} 节的记号。考虑点集
$$
T \subset |\textit{PolarizedCurves}|
$$
其中的点 $\xi$ 满足：存在域 $k$ 以及 $k$ 上表示 $\xi$ 的偶
$(X, \mathcal{L})$，且具有下列两个性质：
\begin{enumerate}
\item $X$ 是亏格为 $g$ 的稳定曲线；
\item $\mathcal{L} = \omega_X^{\otimes 3}$。
\end{enumerate}
显然，在连续映射
$$
|\textit{PolarizedCurves}|
\longrightarrow
|\Curvesstack|
$$
下，集合 $T$ 的像恰为开子集
$$
|\overline{\mathcal{M}}_g| \subset |\Curvesstack|
$$
因此，只需证明 $T$ 拟紧。由引理
\ref{moduli-curves-lemma-polarized-curves-in-polarized}，可知
$$
|\textit{PolarizedCurves}| \subset |\Polarizedstack|
$$
是开闭浸入。因此，只需证明 $T$ 作为 $|\Polarizedstack|$
的子集拟紧。为此，使用《模叠》中的引理
\ref{moduli-lemma-bounded-polarized} 的判据。首先注意，对上述
$(X, \mathcal{L})$，由 Riemann--Roch，Hilbert 多项式 $P$ 是函数
$P(t) = (6g - 6)t + (1 - g)$；见《代数曲线》中的引理
\ref{curves-lemma-rr}。其次，由《代数曲线》中的引理
\ref{curves-lemma-tricanonical}，有 $H^1(X, \mathcal{L}) = 0$，
而且 $\mathcal{L}$ 极丰沛。这恰好意味着：取
% P12-ERR-0313: see ../control/ERRATA_CANDIDATES.jsonl
$n = P(3) - 1$，存在闭浸入
$$
i : X \longrightarrow \mathbf{P}^n_k
$$
使得
% P12-ERR-0314: see ../control/ERRATA_CANDIDATES.jsonl
$\mathcal{L} = i^*\mathcal{O}_{\mathbf{P}^1_k}(1)$，如所求。
\end{proof}

\noindent
下面是本节的主要定理。

\begin{theorem}
\label{moduli-curves-theorem-stable-smooth-proper}
设 $g \geq 2$。代数叠 $\overline{\mathcal{M}}_g$ 是
Deligne--Mumford 叠，在 $\Spec(\mathbf{Z})$ 上真且光滑。
此外，参数化光滑曲线的轨迹 $\mathcal{M}_g$ 是稠密开子叠。
\end{theorem}

\begin{proof}
陈述中提到的大多数性质已经证明。光滑性是引理
\ref{moduli-curves-lemma-stable-curves-smooth}；Deligne--Mumford 性是引理
\ref{moduli-curves-lemma-stable-curves-deligne-mumford}；$\mathcal{M}_g$
的开性是引理 \ref{moduli-curves-lemma-smooth-dense-in-stable}。
由引理 \ref{moduli-curves-lemma-stable-separated}，已知
$\overline{\mathcal{M}}_g \to \Spec(\mathbf{Z})$ 分离；
由引理 \ref{moduli-curves-lemma-stable-quasi-compact}，已知
$\overline{\mathcal{M}}_g$ 拟紧。因此，为证明
$\overline{\mathcal{M}}_g \to \Spec(\mathbf{Z})$ 为真并完成证明，
可以对态射 $\mathcal{M}_g \to \overline{\mathcal{M}}_g$ 和
$\overline{\mathcal{M}}_g \to \Spec(\mathbf{Z})$ 应用
《叠的态射进阶》中的引理
\ref{stacks-more-morphisms-lemma-refined-valuative-criterion-proper}。
因而只需检验以下条件：给定任意 $2$-交换图
$$
\xymatrix{
\Spec(K) \ar[r] \ar[d]_j &
\mathcal{M}_g \ar[r] &
\overline{\mathcal{M}}_g \ar[d] \\
\Spec(A) \ar[rr] & & \Spec(\mathbf{Z})
}
$$
其中 $A$ 是分式域为 $K$ 的离散赋值环，则存在域扩张 $K'/K$
以及支配 $A$ 的赋值环 $A' \subset K'$，使得所诱导的图
$$
\xymatrix{
\Spec(K') \ar[r] \ar[d]_{j'} & \overline{\mathcal{M}}_g \ar[d] \\
\Spec(A') \ar[r] \ar@{..>}[ru] & \Spec(\mathbf{Z})
}
$$
的虚线箭头范畴非空
（《叠的态射》中的定义
\ref{stacks-morphisms-definition-fill-in-diagram}）。
（注意，无需过多担心 $2$-箭头；见《叠的态射》中的引理
\ref{stacks-morphisms-lemma-cat-dotted-arrows-independent}。）
展开其含义，并使用如下事实：$\mathcal{M}_g$ 和
$\overline{\mathcal{M}}_g$ 分别是参数化亏格为 $g$ 的光滑曲线族和
稳定曲线族的代数叠，便发现所需证明的正是稳定约化定理，
即定理 \ref{moduli-curves-theorem-stable-reduction} 所含的结果。
\end{proof}
